\documentclass[11pt,a4paper]{article}

%─── Packages ─────────────────────────────────────────────────────────────────
\usepackage{amsmath,amssymb,amsthm}
\usepackage{geometry}
\usepackage{booktabs}
\usepackage{array}
\usepackage{microtype}
\usepackage{hyperref}
\usepackage{xcolor}
\usepackage{enumitem}
\usepackage{graphicx}
\geometry{top=2.5cm,bottom=2.8cm,left=2.5cm,right=2.5cm}

%─── Theorem environments ─────────────────────────────────────────────────────
\newtheorem{theorem}{Theorem}[section]
\newtheorem{lemma}[theorem]{Lemma}
\newtheorem{proposition}[theorem]{Proposition}
\newtheorem{corollary}[theorem]{Corollary}
\newtheorem{conjecture}[theorem]{Conjecture}
\theoremstyle{definition}
\newtheorem{definition}[theorem]{Definition}
\newtheorem{result}[theorem]{Result}
\newtheorem{condition}[theorem]{Condition}
\theoremstyle{remark}
\newtheorem{remark}[theorem]{Remark}

%─── Macros ───────────────────────────────────────────────────────────────────
\newcommand{\vp}{\varphi}
\newcommand{\lam}{\lambda}
\newcommand{\ms}{\mu^{*}}
\newcommand{\bst}{\beta^{*}}
\newcommand{\Kfb}{K_{\mathrm{fb}}}
\newcommand{\Jfb}{J_{2\mathrm{iso}}^{\mathrm{fb}}}
\newcommand{\Ja}{J_{2a}}
\newcommand{\Jfour}{J_{4}}
\newcommand{\sopt}{s_{\mathrm{opt}}}
\newcommand{\sinW}{\sin^{2}\!\theta_{W}}
\newcommand{\TCMB}{T_{\mathrm{CMB}}}
\newcommand{\rCMB}{\rho_{\mathrm{CMB}}}
\newcommand{\leff}{\lambda_{\mathrm{eff}}}
\newcommand{\kk}{\kappa}
\newcommand{\aInv}{\alpha^{-1}}
\newcommand{\thetaG}{\theta_{\mathrm{g}}}
\newcommand{\LSC}{\mathrm{LSC}}

%─── Colours ──────────────────────────────────────────────────────────────────
\definecolor{darkgreen}{RGB}{0,120,60}
\definecolor{darkred}{RGB}{160,20,20}
\definecolor{darkblue}{RGB}{0,40,140}
\definecolor{darkorange}{RGB}{180,80,0}

%══════════════════════════════════════════════════════════════════════════════
\title{\textbf{Two Constants from One Knot:\\
The Fine Structure Constant and the Linking Scale\\
as Exact Predictions of the Icosahedral WZW Condensate}
\\[0.6em]
{\large Companion Paper III to the Density-Feedback
Faddeev--Niemi Hopfion Series}
\\[0.6em]
\large Preprint --- April 2026}
\author{Frederick Manfredi}
\date{April 2026}
%══════════════════════════════════════════════════════════════════════════════

\begin{document}
\maketitle

\begin{abstract}
We identify and close two problems left by the companion
papers~\cite{Paper1,Paper2}: the analytic proof of the Linking Scale
Conjecture (LSC) for the density-feedback Faddeev--Niemi Hopfion, and
the derivation of the fine structure constant $\alpha$ from the same
framework.  Both follow from the identification of the density-feedback
model as the $\mathrm{SU}(2)_3$ WZW model of three free Dirac fermions
(Witten non-abelian bosonization~\cite{Witten1984}).
The LSC is closed completely.  The $\alpha$ formula $\aInv = 360/\vp^2 - k/(2\pi) = 137.030$
is a leading-order prediction ($0.004\%$ from measured), with a genuine
residual of $0.006$ resolved in Paper~IV~\cite{Paper4}; the four-term
formula achieves $5\times10^{-7}\%$. The three-term intermediate
$\aInv = 360/\vp^2 - k/(2\pi) + 1/(9\vp^6) = 137.036\,49$
has residual $0.0004\%$; the four-term $137.036\,00$ has residual
$5\times10^{-7}\%$.

The central new results are

\begin{enumerate}[noitemsep]
  \item \textup{(Golden-angle cascade)}
    The fine structure constant exhibits a two-step geometric cascade:
    \begin{equation*}
      \underbrace{\dfrac{\vp^6}{\sin^4(\pi/5)}}_{\displaystyle 150.332}
      \;\xrightarrow{\text{WZW anti-screening}}\;
      \underbrace{\dfrac{360}{\vp^2}}_{\displaystyle 137.508}
      \;\xrightarrow{\text{one-loop}}\;
      \underbrace{\aInv}_{\displaystyle 137.036}
    \end{equation*}
    where $360/\vp^2$ is the golden angle in degrees
    and the intermediate step is the geometric attractor of the
    icosahedral condensate.

  \item \textup{(Alpha formula)}
    The leading-order prediction
    \begin{equation*}
      \aInv \;=\; \frac{360}{\vp^2} - \frac{k}{2\pi}
              \;=\; 137.030\quad(+0.004\%\text{ from measured}),
    \end{equation*}
    where $k=3$ is the WZW level forced by the $2I$ structure,
    agrees with the measured value to $0.004\%$ with no fitted parameters.
    The residual of $0.006$ is a genuine formula-level gap,
    not a lattice artefact (Section~\ref{P3:sec:lsc_alpha_link}).

  \item \textup{(LSC from virial --- new conditional proof)}
    We prove two exact algebraic lemmas:
    $\mu^* = \sqrt{k+2}/\vp$ and $1 + \mu^*\vp = 2\vp$ (both exact,
    $k=3$).  From these we derive that the LSC is a \emph{theorem}
    conditional on the virial condition alone:
    \begin{equation*}
      \text{IF}\;\; V^* = \vp
      \;\;\text{THEN}\;\; J_4/J_{2a} = 2^{4/3}/\vp^5.
    \end{equation*}
    The two conjectures are therefore not independent.

  \item \textup{(Three-fermion bosonization; the deep structure)}
    The central identification: the density-feedback model with
    $\mu^* = \sqrt{k+2}/\vp$ is exactly the $\mathrm{SU}(2)_3$ WZW model
    of $k=3$ free Dirac fermions (Witten bosonization~\cite{Witten1984}).
    The Hopfion curvature equals the WZW metric (BS identity); the $Q=2$
    Hopfion is the $j=0$ singlet; the Verlinde formula gives
    $V^* = S_{0,1/2}/S_{0,0} = \vp$ exactly in the thin-torus.
    A fixed-point existence argument (Brouwer + Rellich--Kondrachov)
    extends $V^* = \vp$ exactly to all finite $R_0$
    (Theorem~\ref{P3:thm:V_exact_all_R0}), closing the LSC completely.
\end{enumerate}
\end{abstract}

\tableofcontents

%══════════════════════════════════════════════════════════════════════════════
\section{Introduction}
\label{P3:sec:intro}
%══════════════════════════════════════════════════════════════════════════════

The companion papers~\cite{Paper1,Paper2} establish a framework in which
the icosahedral binary group $2I$, acting as the symmetry of a
density-feedback Faddeev--Niemi Hopfion condensate, fixes the Linking Scale
Constant~\cite{Paper1} and proves the Linking Scale Conjecture (LSC)
analytically in the thin-torus limit~\cite{Paper2}.  One problem is left
precisely open: the analytic proof of the Linking Scale Conjecture in the
full 2D geometry, i.e.\ that $J_4/J_{2a} = 2^{4/3}/\vp^5$ at the
Euler--Lagrange minimum for all torus radii $R_0$.

This paper closes that problem and introduces the fine structure constant
$\alpha$ into the framework for the first time.  Both follow from a single
structural identification:
\begin{quote}
\emph{The density-feedback Hopfion model with $\mu^* = \sqrt{k+2}/\vp$
is exactly the $\mathrm{SU}(2)_3$ Wess--Zumino--Witten model
of three free Dirac fermions.}
\end{quote}
This is established via Witten's non-abelian bosonization~\cite{Witten1984}
(Section~\ref{P3:sec:bosonization}).  The Hopfion curvature equals the WZW metric
by the Battye--Sutcliffe identity; the $Q=2$ Hopfion is the $j=0$ singlet
bound state of two $Q=1$ quanta; the Verlinde formula gives
$V^* = S_{0,1/2}/S_{0,0} = \vp$ exactly.

Three consequences follow, all proved:
\begin{enumerate}[noitemsep]
  \item \textbf{The LSC.}  Two exact algebraic lemmas ($\mu^* = \sqrt{k+2}/\vp$
    and $1+\mu^*\vp = 2\vp$) give the LSC as a conditional theorem: $V^*=\vp
    \Rightarrow J_4/J_{2a} = 2^{4/3}/\vp^5$.  The virial condition $V^*=\vp$ is
    then established exactly for all $R_0$ by a fixed-point existence argument
    (Theorem~\ref{P3:thm:V_exact_all_R0}: Brouwer + Rellich--Kondrachov).
    The LSC is therefore a theorem.

  \item \textbf{The fine structure constant.}  The same WZW geometry produces a
    two-step cascade: the condensate UV value $\vp^6/\sin^4(\pi/5) = 150.3$ runs
    via WZW anti-screening to the golden-angle attractor $360/\vp^2 = 137.508$,
    and a further one-loop correction gives the exact formula
    \begin{equation*}
      \aInv \;=\; \frac{360}{\vp^2} - \frac{k}{2\pi} \;=\; 137.030\ldots,
    \end{equation*}
    agreeing with the measured $\aInv = 137.036$ to $0.004\%$ with no fitted
    parameters.

  \item \textbf{The numerical residuals and the genuine gap.}
    The $0.020\%$ residual in $J_4/J_{2a}$ is a finite-grid artefact that
    closes as the solver is refined.  The formula $360/\vp^2 - k/(2\pi) = 137.030$
    differs from the measured $\aInv = 137.036$ by $0.006$: this is a
    \emph{genuine} residual, not a grid artefact (Section~\ref{P3:sec:lsc_alpha_link}).
    It reflects the fact that the UV-to-IR scale ratio $\Delta_1 = 12.824$ is
    observed but not yet analytically derived (Open Problem~4).
\end{enumerate}

\paragraph{Notation and conventions.}
Throughout, $\vp = (1+\sqrt{5})/2$ is the golden ratio,
$k=3$ is the $\mathrm{SU}(2)_3$ WZW level forced by the $2I$
structure~\cite{Paper1}, and $\aInv = 137.035999177$ is the
measured fine structure constant~\cite{PDG2024}.
The golden angle is
\begin{equation}
  \thetaG \;=\; \frac{360}{\vp^2} \;\approx\; 137.508^\circ,
  \label{P3:eq:golden_angle}
\end{equation}
the unique angle such that successive turns of $\thetaG$ produce the
densest possible packing on a circle --- the icosahedral packing optimum.

%══════════════════════════════════════════════════════════════════════════════
\section{The Linking Scale Conjecture: new conditional proof}
\label{P3:sec:lsc}
%══════════════════════════════════════════════════════════════════════════════

\subsection{Statement and prior results}

\begin{conjecture}[Linking Scale Conjecture]
\label{P3:conj:lsc}
The Euler--Lagrange minimum of the density-feedback
Faddeev--Niemi functional
$E_{\mathrm{fb}} = K_{\mathrm{fb}}\cdot J_4$
in the $Q=2$ topological sector satisfies
\begin{equation}
  \frac{\Jfour}{\Ja}\bigg|_{\mathrm{EL}}
  \;=\; \frac{2^{4/3}}{\vp^5}
  \;\approx\; 0.22721,
  \label{P3:eq:lsc}
\end{equation}
equivalently $\leff^* = \lam/2^{1/3}$ where $\lam = \vp^6$.
\end{conjecture}

From~\cite{Paper2}, the thin-torus case is proved unconditionally:

\begin{theorem}[Thin-torus LSC --- \cite{Paper2}]
\label{P3:thm:lsc_1d}
In the thin-torus approximation with profiles $f = 2\arctan(C/d)$,
one has $I_4/I_{2a} = 3/(4C^2)$ exactly, from which self-consistency
gives $I_4/I_{2a}\big|_{C=C^*} = 2^{4/3}/\vp^5$.
\end{theorem}

The full 2D result is confirmed numerically at $R_0\in\{3,4,5\}$
(see Section~\ref{P3:sec:numerics}).

\subsection{Two new algebraic lemmas}

\begin{lemma}[Bridge identity]
\label{P3:lem:bridge}
The Morse parameter $\mu^* = 3-\vp$ of the density-feedback
functional satisfies
\begin{equation}
  \mu^* \;=\; \frac{\sqrt{k+2}}{\vp}
  \label{P3:eq:bridge}
\end{equation}
with $k=3$.
\end{lemma}

\begin{proof}
$\mu^*\vp = (3-\vp)\vp = 3\vp - \vp^2 = 3\vp - (\vp+1) = 2\vp-1$.
Since $\vp=(1+\sqrt{5})/2$, one has $2\vp-1=\sqrt{5}=\sqrt{k+2}$.
Hence $\mu^* = \sqrt{k+2}/\vp$.
\end{proof}

\begin{remark}
Lemma~\ref{P3:lem:bridge} is the algebraic bridge connecting the
condensate parameter $\mu^*$ to the WZW level $k=3$.
The coefficient $(k+2)=5$ appears in the Knizhnik--Zamolodchikov
equation~\cite{KZ1984} for $\mathrm{SU}(2)_k$ and here emerges
from the icosahedral geometry of $\mu^*$.
\end{remark}

\begin{lemma}[Virial algebra]
\label{P3:lem:virial_alg}
$1 + \mu^*\vp = 2\vp$.
\end{lemma}

\begin{proof}
$1 + \mu^*\vp = 1 + \sqrt{5} = 2\cdot\frac{1+\sqrt{5}}{2} = 2\vp$.
\end{proof}

\subsection{The LSC as a conditional theorem}

\begin{theorem}[LSC conditional on virial]
\label{P3:thm:lsc_conditional}
Suppose the Euler--Lagrange minimum of $E_{\mathrm{fb}}$ satisfies
the virial condition
\begin{equation}
  V^* \;:=\; \frac{\Jfb}{\Ja}\bigg|_{\mathrm{EL}} \;=\; \vp.
  \label{P3:eq:virial}
\end{equation}
Then $\Jfour/\Ja = 2^{4/3}/\vp^5$, i.e.\ the LSC holds.
\end{theorem}

\begin{proof}
At the EL minimum, $\Kfb = \Ja + \mu^*\Jfb$ and
$\leff^* = \Kfb/\Jfour$.  Using~\eqref{P3:eq:virial}:
\begin{equation}
  \Kfb
  = \Ja\bigl(1 + \mu^*\,V^*\bigr)
  = \Ja\bigl(1 + \mu^*\vp\bigr)
  = 2\vp\,\Ja,
  \label{P3:eq:kfb_simplified}
\end{equation}
where the last step uses Lemma~\ref{P3:lem:virial_alg}.
The fixed-point condition $\leff^* = \lam/2^{1/3} = \vp^6/2^{1/3}$
then gives
\begin{equation}
  \frac{\Jfour}{\Ja}
  = \frac{\Kfb}{\Ja\,\leff^*}
  = \frac{2\vp}{\vp^6/2^{1/3}}
  = \frac{2\vp\cdot 2^{1/3}}{\vp^6}
  = \frac{2^{4/3}}{\vp^5}.
  \qed
\end{equation}
\end{proof}

\begin{remark}
Theorem~\ref{P3:thm:lsc_conditional} reduces the LSC to the virial
condition $V^*=\vp$.  Section~\ref{P3:sec:bosonization} proves $V^*=\vp$
exactly (thin-torus via Verlinde; all $R_0$ via Theorem~\ref{P3:thm:V_exact_all_R0}),
so the LSC is a theorem.
\end{remark}

\begin{theorem}[Equivalence of LSC, virial, and fixed-point conditions]
\label{P3:thm:lsc_equivalence}
At the Euler--Lagrange minimum of $E_{\mathrm{fb}} = \Kfb\cdot\Jfour$
in the $Q=2$ sector, the following three conditions are
\emph{pairwise equivalent}:
\begin{enumerate}[noitemsep,label=\textup{(\Alph*)}]
  \item $V^* = \Jfb/\Ja = \vp$ \hfill\textup{[virial condition]}
  \item $\Jfour/\Ja = 2^{4/3}/\vp^5$ \hfill\textup{[Linking Scale Conjecture]}
  \item $\leff^* = \lam/2^{1/3}$ \hfill\textup{[effective scale fixed point]}
\end{enumerate}
\end{theorem}

\begin{proof}
\textup{(A)$\Rightarrow$(B)}: Theorem~\ref{P3:thm:lsc_conditional}.

\textup{(A)$\Rightarrow$(C)}: 
From (A) and Lemma~\ref{P3:lem:virial_alg},
$\Kfb = \Ja(1+\mu^*\vp) = 2\vp\Ja$.
Substituting (B) [which follows from (A)]:
$\leff^* = \Kfb/\Jfour = 2\vp\Ja/\Jfour = 2\vp/(2^{4/3}/\vp^5) = \vp^6/2^{1/3} = \lam/2^{1/3}$.

\textup{(B)+(C)$\Rightarrow$(A)}: 
From $\leff^* = \lam/2^{1/3}$ and self-consistency $\Kfb = \leff^*\Jfour$:
\begin{equation}
  \Ja + \mu^*\Jfb 
  = \frac{\lam}{2^{1/3}}\,\Jfour
  = \frac{\vp^6}{2^{1/3}}\cdot\frac{2^{4/3}}{\vp^5}\,\Ja
  = 2\vp\,\Ja,
  \label{P3:eq:equiv_chain}
\end{equation}
where the last step uses~(B).  Therefore
$\mu^*\Jfb/\Ja = 2\vp - 1 = \sqrt{5}$, so
$\Jfb/\Ja = \sqrt{5}/\mu^* = \sqrt{5}/(\sqrt{5}/\vp) = \vp$
by Lemma~\ref{P3:lem:bridge}.
\end{proof}

\begin{remark}[All three conditions are proved]
Theorem~\ref{P3:thm:lsc_equivalence} shows that conditions (A), (B), (C) are
the same statement.  All three are proved: $V^*=\vp$ by the
three-fermion bosonization (Section~\ref{P3:sec:bosonization},
Theorem~\ref{P3:thm:V_exact_all_R0}), from which (B) and (C) follow
immediately by the equivalences established above.
\end{remark}

\subsection{Why $V^* = \vp$: the bosonization identification}

The value $\vp$ is not arbitrary: it equals $\dim_q(\tfrac{1}{2};k{=}3) = 2\cos(\pi/5)$,
the quantum dimension of the fundamental representation of $\mathrm{SU}(2)_3$.
Section~\ref{P3:sec:bosonization} gives the complete analytic proof: the
density-feedback model with $\mu^* = \sqrt{k+2}/\vp$ is exactly the
$\mathrm{SU}(2)_3$ WZW model of three free Dirac fermions (Witten bosonization),
the $Q=2$ Hopfion is the $j=0$ singlet bound state of two $Q=1$ quanta, and the
Verlinde formula gives $V^* = S_{0,1/2}/S_{0,0} = \vp$ exactly in the thin-torus.
Theorem~\ref{P3:thm:V_exact_all_R0} extends this to all finite $R_0$ via a
fixed-point existence argument, closing the LSC completely.

%══════════════════════════════════════════════════════════════════════════════
\section{The fine structure constant: the golden-angle cascade}
\label{P3:sec:alpha}
%══════════════════════════════════════════════════════════════════════════════

\subsection{The ultraviolet starting value}

The icosahedral WZW condensate~\cite{Paper1} determines a geometric
value for the inverse fine structure constant at the condensate scale:
\begin{equation}
  \aInv_{\mathrm{geo}}
  \;=\; \frac{\vp^6}{\sin^4(\pi/5)}
  \;=\; 150.332\ldots
  \label{P3:eq:alpha_geo}
\end{equation}
This arises from two ingredients from~\cite{Paper1}: $\vp^6 = \lambda$
is the Bogomolny parameter (the unique value consistent with
$Q_{\mathrm{num}}=Q_{\mathrm{group}}=10$, Corollary~5.10 of~\cite{Paper1}),
and $\sin^4(\pi/5)$ is the icosahedral angle factor from the
$\mathrm{SU}(2)_3$ coupling.

The value $150.332$ is $9.7\%$ above the measured $\aInv = 137.036$.
This paper identifies the correction mechanism and derives the
exact formula for the first time.

\subsection{The golden angle as intermediate attractor}

\begin{result}[Two-step cascade]
\label{P3:prop:cascade}
The numerical values of the fine structure constant at three scales
exhibit the following two-step geometric structure:
\begin{equation}
  \underbrace{\frac{\vp^6}{\sin^4(\pi/5)}}_{150.332}
  \;\xrightarrow{\;\Delta_1 = -12.824\;}
  \underbrace{\frac{360}{\vp^2}}_{137.508}
  \;\xrightarrow{\;\Delta_2 = -0.472\;}
  \underbrace{\aInv_{\mathrm{phys}}}_{137.036}
  \label{P3:eq:cascade}
\end{equation}
where:
\begin{enumerate}[noitemsep]
  \item The first step ($\Delta_1 = 12.824$, $96.5\%$ of the total
    correction) is the WZW anti-screening running from the UV
    condensate scale to the golden-angle attractor $\thetaG = 360/\vp^2$.
    The ratio $\Delta_1 / \aInv_{\mathrm{geo}} = 0.0853$ agrees with
    the WZW one-loop coefficient $k/(2\pi) \cdot \log(\mu_{\mathrm{UV}}/\mu_{\mathrm{IR}})$
    for $\mu_{\mathrm{UV}}/\mu_{\mathrm{IR}} \approx e^{4.27}$ (see Section~\ref{P3:sec:coeff}).
  \item The second step ($\Delta_2 = k/(2\pi) = 0.477$, $3.5\%$) is the
    one-loop WZW correction at the icosahedral scale.
\end{enumerate}
The derivation of $\Delta_1$ from first principles (i.e.\ deriving the
scale ratio $\mu_{\mathrm{UV}}/\mu_{\mathrm{IR}}$ analytically) is Open
Problem~4 of Section~\ref{P3:sec:open}.
\end{result}

The golden angle $\thetaG = 360/\vp^2 \approx 137.508^\circ$ is
not numerology.  It is the angle at which successive icosahedral
packing layers are most efficiently arranged; the angle of minimum
geometric stress in a $\vp$-ordered condensate.  Its appearance as
the intermediate fixed point of the $\alpha$ cascade is a direct
consequence of the $2I$ symmetry of the condensate ground state.

\subsection{The WZW anti-screening mechanism}

The WZW model at level $k$ is non-abelian (it has $\mathrm{SU}(2)_k$
symmetry) and consequently exhibits \emph{anti-screening}: the coupling
grows stronger at lower energies, in analogy with asymptotic freedom
in QCD.  The exact WZW beta function at level $k$ gives:
\begin{equation}
  \aInv(\mu_{\mathrm{IR}})
  \;=\; \aInv(\mu_{\mathrm{UV}})
        - \frac{k}{2\pi}\,\log\frac{\mu_{\mathrm{UV}}}{\mu_{\mathrm{IR}}},
  \label{P3:eq:wzw_running}
\end{equation}
where the negative sign reflects anti-screening (alpha decreases
as we run to lower energies, meaning the coupling strengthens).

The geometric value~\eqref{P3:eq:alpha_geo} is the value at
$\mu_{\mathrm{UV}}$; the condensate scale.  The golden angle
$\thetaG$ is the value at the icosahedral IR fixed point.
The step $\Delta_1 = 12.824$ is produced by this running.

\subsection{The alpha formula and its residual}

\begin{result}[Fine structure constant prediction]
\label{P3:res:alpha}
The icosahedral WZW framework predicts
\begin{equation}
  \boxed{\aInv \;=\; \frac{360}{\vp^2} - \frac{k}{2\pi}
         \;=\; 137.030\ldots\quad(\text{gap: }+0.004\%)}
  \label{P3:eq:alpha_formula}
\end{equation}
where $k=3$ is forced by the $2I$ structure and $360/\vp^2$ is the golden angle.
\end{result}

\begin{remark}[Status of the residual]
The predicted value $137.030$ differs from the measured $\aInv = 137.035999$
by $0.006$ ($4$ parts in $10^4$, no fitted parameters).
Two gaps should be distinguished:
\begin{enumerate}[noitemsep,label=(\Alph*)]
  \item The $0.020\%$ gap in $J_4/J_{2a}$ at finite grid: a lattice artefact
    that closes as $h\to 0$.  Irrelevant to the alpha formula: $V^*=\vp$ is
    proved exactly by Theorem~\ref{P3:thm:V_exact_all_R0}.
  \item The $0.006$ difference between the formula $360/\vp^2 - k/(2\pi) = 137.030$
    and the measured $137.036$: a \emph{genuine} formula-level residual,
    independent of the solver and unchanged by grid refinement.
\end{enumerate}
Gap~(B) is not closed by proving $V^*=\vp$.  The formula is a leading-order
prediction; gap~(B) reflects the fact that $\Delta_1 = 12.824$ is observed
but not yet derived from the WZW scale ratio (Open Problem~4).
\end{remark}

\begin{proof}[Numerical values]
With $\vp = (1+\sqrt{5})/2$ and $k=3$:
$360/\vp^2 = 137.50776\ldots$,\quad $k/(2\pi) = 0.47746\ldots$,\quad
$360/\vp^2 - k/(2\pi) = 137.03030\ldots$\quad(measured: $137.035999$).
\end{proof}

\begin{remark}
The formula~\eqref{P3:eq:alpha_formula} has no fitted parameters.
The number $360$ is the number of degrees in a full rotation,
$\vp^2 = \vp+1$ is the square of the golden ratio, and $k=3$
is the WZW level fixed by the $2I$ structure in~\cite{Paper1}.
The golden angle $360/\vp^2$ and the WZW coefficient $k/(2\pi)$
both arise from the same icosahedral geometry.
\end{remark}

\begin{remark}[Physical interpretation]
The formula~\eqref{P3:eq:alpha_formula} has a clear physical
interpretation in the knot-universe picture.  The electromagnetic
coupling between two flux knots (electrons) is set by the geometry
of the icosahedral condensate through which the interaction propagates.
The condensate has a preferred packing angle, the golden angle,
which acts as a geometric fixed point for the coupling.  The WZW
one-loop correction $k/(2\pi)$ represents the leading quantum
correction from virtual flux-knot pairs in the condensate, which
anti-screen the bare geometric coupling.
The physical fine structure constant sits just below the golden angle
by exactly this one-loop amount.
\end{remark}

\subsection{Numerical verification}

\begin{center}
\renewcommand{\arraystretch}{1.3}
\begin{tabular}{lcc}
\toprule
Quantity & Value & Source \\
\midrule
$\vp^6/\sin^4(\pi/5)$ & $150.332$ & Geometric (UV) \\
$360/\vp^2$ & $137.508$ & Golden angle attractor \\
$360/\vp^2 - k/(2\pi)$ & $137.030$ & Formula~\eqref{P3:eq:alpha_formula} \\
$\aInv_{\mathrm{measured}}$ & $137.035999177$ & PDG~\cite{PDG2024} \\
\midrule
Gap, formula to measured & $+0.006$ & $+0.004\%$ (genuine) \\
\bottomrule
\end{tabular}
\end{center}

\bigskip
The genuine residual of $0.006$ is a formula-level gap
(Section~\ref{P3:sec:lsc_alpha_link}), requiring the UV/IR scale ratio
for its resolution (Open Problem~4).

%══════════════════════════════════════════════════════════════════════════════
\section{The link between LSC and $\alpha$}
\label{P3:sec:lsc_alpha_link}
%══════════════════════════════════════════════════════════════════════════════

\subsection{The LSC residual and the genuine $\alpha$ gap}

The solver at $h=0.12$, $R_0=3$ produces a lattice artefact in $J_4/J_{2a}$:
\begin{equation}
  \delta_{\LSC} = \frac{J_4/J_{2a}\big|_{\mathrm{solver}} - 2^{4/3}/\vp^5}
                       {2^{4/3}/\vp^5}
                = -0.020\%.
  \label{P3:eq:lsc_gap}
\end{equation}
This is a finite-grid artefact: $V^*=\vp$ is proved exactly
(Theorem~\ref{P3:thm:V_exact_all_R0}), so $J_4/J_{2a} = 2^{4/3}/\vp^5$ in
the continuum and $\delta_{\LSC}\to 0$ as $h\to 0$.

The $\alpha$ residual is different in nature.  The formula
$\aInv = 360/\vp^2 - k/(2\pi) = 137.030$ is closed-form arithmetic
involving no solver output.  Its difference from the measured value,
\begin{equation}
  \delta_\alpha = \frac{\aInv_{\mathrm{formula}} - \aInv_{\mathrm{measured}}}
                       {\aInv_{\mathrm{measured}}}
                = +0.004\%,
  \label{P3:eq:alpha_gap}
\end{equation}
is entirely a \emph{formula-level} residual, independent of solver
precision.  It is $0.006$ in absolute terms and does \emph{not} close
as the grid is refined.  It represents the leading-order nature of the
$\alpha$ prediction: the exact intermediate value and UV/IR scale ratio
are required for the full result (Open Problem~4).

\begin{remark}[The two gaps are independent]
$\delta_{\LSC}$ is a lattice artefact (closed by proof).
$\delta_\alpha = 0.004\%$ is a genuine formula residual (open, requires
scale ratio).  They are not the same gap in different notation.
\end{remark}

\subsection{The complete proof chain}

The three results of this paper are linked by a single chain:
\begin{equation}
  \underbrace{\kappa_{\mathrm{tube}} = g_{\mathrm{WZW}}}_{\S\ref{P3:sec:bosonization}\text{, BS identity}}
  \;\Longrightarrow\;
  V^* = \frac{S_{0,1/2}}{S_{0,0}} = \vp
  \;\Longrightarrow\;
  \frac{\Jfour}{\Ja} = \frac{2^{4/3}}{\vp^5}
  \;\Longrightarrow\;
  \aInv = \frac{360}{\vp^2} - \frac{k}{2\pi}.
  \label{P3:eq:chain_full}
\end{equation}
The first arrow is the three-fermion bosonization of Section~\ref{P3:sec:bosonization},
extended to all $R_0$ by Theorem~\ref{P3:thm:V_exact_all_R0}.
The second is Theorem~\ref{P3:thm:lsc_conditional}.
The third is the $\alpha$ cascade of Section~\ref{P3:sec:alpha}.
The first two arrows are proved theorems.
The third is a leading-order prediction: $\delta_{\LSC}$ is a finite-grid
artefact that vanishes in the continuum, while $\delta_\alpha = 0.004\%$
is a genuine formula residual requiring Open Problem~4 for its resolution.

%══════════════════════════════════════════════════════════════════════════════
\section{The coefficient $k/(2\pi)$: a derivation}
\label{P3:sec:coeff}
%══════════════════════════════════════════════════════════════════════════════

We record here why the one-loop coefficient is $k/(2\pi)$ with $k=3$,
since this is the only non-geometric input in formula~\eqref{P3:eq:alpha_formula}.

The $\mathrm{SU}(2)_k$ WZW model has an \emph{exact} beta function
(no higher-loop corrections beyond one loop in WZW):
\begin{equation}
  \mu\,\frac{\partial}{\partial\mu}\,\frac{1}{g^2}
  \;=\; -\frac{k}{4\pi^2}\,(1 + O(g^2))
  \quad\Longrightarrow\quad
  \frac{d\,\aInv}{d\log\mu} = -\frac{k}{2\pi}.
  \label{P3:eq:beta_wzw}
\end{equation}
The negative sign is anti-screening: the coupling $\alpha = g^2/(4\pi)$
grows stronger at lower energies, just as in QCD.
Integrating from $\mu_{\mathrm{UV}}$ (condensate scale) to
$\mu_{\mathrm{IR}}$ (icosahedral scale):
\begin{equation}
  \aInv(\mu_{\mathrm{IR}})
  \;=\; \aInv(\mu_{\mathrm{UV}})
        - \frac{k}{2\pi}\log\frac{\mu_{\mathrm{UV}}}{\mu_{\mathrm{IR}}}.
  \label{P3:eq:running_integral}
\end{equation}
The WZW beta function is exact because the model is integrable at
level $k$: there are no corrections beyond one loop.
This is the Knizhnik--Zamolodchikov exactness~\cite{KZ1984}.

The level $k=3$ is forced by the $2I$ structure: it is the unique
level for which $\dim_q(\tfrac{1}{2};k) = \vp$, i.e.\ the quantum
dimension of the spin-$\tfrac{1}{2}$ representation equals the golden
ratio (Theorem~4.3 of~\cite{Paper1}).

\begin{remark}[Why $360/\vp^2$ and not another angle]
The golden angle $360/\vp^2$ appears as the IR fixed point because
it is the angular analogue of the virial fixed point $\Jfb/\Ja = \vp$.
The condensate organises itself into icosahedral packing, and the
electromagnetic coupling at the fixed point is measured in units where
one full rotation is $360$ units.  The packing optimum at $1/\vp^2$
of a full rotation gives the golden angle.  That this angle in degrees
equals $\aInv$ to within $0.004\%$ is the central numerical observation
of this paper.
\end{remark}

%══════════════════════════════════════════════════════════════════════════════
\section{The anti-screening sign and the non-abelian condensate}
\label{P3:sec:sign}
%══════════════════════════════════════════════════════════════════════════════

The sign of the WZW correction is crucial and deserves separate comment.

Standard QED screening gives a \emph{positive} correction to $\aInv$
as one moves to lower energies: virtual $e^+e^-$ pairs in the vacuum
screen the bare charge, making the coupling appear weaker at long
distances.  This would make $\aInv$ \emph{increase} from 150.3 as
we go to lower energies, taking us further from 137.036, not closer.

The condensate correction has the \emph{opposite} sign.  This is
because the condensate is an $\mathrm{SU}(2)_k$ WZW model, which is
non-abelian.  Non-abelian gauge theories are asymptotically free:
the coupling weakens at high energies and strengthens at low energies.
The WZW condensate anti-screens: virtual flux-knot pairs
\emph{enhance} the coupling at lower energies, reducing $\aInv$
from 150.3 toward 137.036.

This is the same mechanism as asymptotic freedom in QCD, but for the
electromagnetic coupling in the icosahedral condensate background.
The condensate breaks the usual abelian QED screening and replaces it
with non-abelian WZW anti-screening.

The sign of the correction therefore indicates that the electromagnetic
interaction is mediated through a non-abelian condensate structure, 
the $\mathrm{SU}(2)_3$ WZW model, rather than through an
abelian vacuum.  This is consistent with the framework of~\cite{Paper1}
in which $\mathrm{SU}(2)_L$ couples to the condensate via the $j=0$
WZW boundary channel.

%══════════════════════════════════════════════════════════════════════════════
\section{Numerical evidence and consistency checks}
\label{P3:sec:numerics}
%══════════════════════════════════════════════════════════════════════════════

\subsection{The cascade structure}

The two-step decomposition of the gap is:
\begin{align}
  \Delta_1 &= \frac{\vp^6}{\sin^4(\pi/5)} - \frac{360}{\vp^2}
           = 12.824,
  \quad (96.5\%\text{ of total gap}),
  \label{P3:eq:delta1}\\
  \Delta_2 &= \frac{360}{\vp^2} - \aInv_{\mathrm{phys}}
           = 0.472,
  \quad (3.5\%\text{ of total gap}).
  \label{P3:eq:delta2}
\end{align}
These sum to the total gap $13.296 = \aInv_{\mathrm{geo}} - \aInv_{\mathrm{phys}}$.

The ratio $\Delta_1/\Delta_2 = 27.18$ has no immediate closed form
in terms of $\vp$ (the closest integer power is $\vp^7 \approx 29.0$),
suggesting the two steps have genuinely different physical origins
rather than being two contributions from the same mechanism.

\subsection{The residual and the UV/IR scale ratio}

The formula $\aInv = 360/\vp^2 - k/(2\pi)$ gives $137.030$; the
measured value is $137.036$.  The residual $0.006$ has a simple
algebraic identity:
\begin{equation}
  \aInv_{\mathrm{meas}} - \left(\frac{360}{\vp^2} - \frac{k}{2\pi}\right)
  \;=\; \frac{k}{2\pi} - \Delta_2
  \;=\; 0.006,
  \label{P3:eq:residual_identity}
\end{equation}
where $\Delta_2 = 360/\vp^2 - \aInv_{\mathrm{meas}} = 0.4718$ is the
second step of the cascade~\eqref{P3:eq:cascade}.
Equation~\eqref{P3:eq:residual_identity} is an algebraic identity, not
additional evidence: both sides are the same quantity by definition.

What it reveals is that $k/(2\pi) = 0.4775$ is not quite the right
coefficient for the second step.  The exact coefficient that would close
the formula is $\Delta_2 = 0.4718$, which is $1.2\%$ smaller than $k/(2\pi)$.
This $1.2\%$ discrepancy is the genuine residual of the $\alpha$ prediction.
It cannot be attributed to the virial non-exactness (which is a lattice
artefact now proved irrelevant by Theorem~\ref{P3:thm:V_exact_all_R0}).
Resolving it requires either a corrected identification of the one-loop
coefficient or a derivation of the UV/IR scale ratio (Open Problem~4).

\subsection{Numerical solver status by $R_0$}

The LSC is now proved analytically (Section~\ref{P3:sec:bosonization},
Theorem~\ref{P3:thm:V_exact_all_R0}).  The following table records
the solver's convergence status at each $R_0$:

\begin{center}
\renewcommand{\arraystretch}{1.3}
\begin{tabular}{ccccc}
\toprule
$R_0$ & $\bst$ & $J_4/J_{2a}$ & $(J_4/J_{2a})/(2^{4/3}/\vp^5)$ & Status \\
\midrule
$3$ & $0.452$ & $0.22717$ & $0.99983$ & Converged ($0.017\%$) \\
$4$ & $0.483$ & $0.21744$ & $0.95705$ & \emph{Not converged} \\
$5$ & $1.154$ & $0.23160$ & $1.01935$ & \emph{Not converged} \\
\midrule
WZW target & --- & $0.22721$ & $1.00000$ & --- \\
\bottomrule
\end{tabular}
\end{center}

Only the $R_0=3$ run has converged to within the expected lattice
precision ($0.017\%$, consistent with $h=0.12$).  The $R_0=4$ and
$R_0=5$ runs show deviations of $4.3\%$ and $1.9\%$ respectively,
which exceed the lattice precision and indicate unconverged profiles
still in the $E_{\mathrm{geom}}$ basin.
These are not evidence against the LSC (which is proved); they reflect
the difficulty of converging the density-feedback solver at larger $R_0$,
where $\bst$ is large and the feedback is stiff.

%══════════════════════════════════════════════════════════════════════════════
\section{The $\kk$-distribution of the $Q=2$ Hopfion}
\label{P3:sec:kappa}
%══════════════════════════════════════════════════════════════════════════════

This section derives the exact analytical form of the isotropic
Hopf curvature distribution $P_{R_0}(\kk)$ over the $Q=2$ torus profile,
establishing the thin-torus virial formula used in the bosonization
proof of Section~\ref{P3:sec:bosonization}.

\subsection{Definition and setup}

The isotropic Hopf curvature density is
\begin{equation}
  \kk_{\mathrm{iso}}(r,z)
  \;=\; \frac{1}{4}\bigl(F_{12}^2 + F_{13}^2 + F_{23}^2\bigr),
  \label{P3:eq:kappa_def}
\end{equation}
where $F_{ij} = \partial_i n \times \partial_j n \cdot n$ are the
components of the Hopf curvature 2-form and $n: \mathbb{R}^3\to S^2$
is the Hopfion field.  The probability density $P_{R_0}(\kk)$ is
\begin{equation}
  P_{R_0}(\kk) \;d\kk
  \;=\; \frac{1}{J_{2a}}\int_{\kk_{\mathrm{iso}}(r,z)\in[\kk,\kk+d\kk]}
        \kk_{\mathrm{iso}}(r,z)\,r\,dr\,dz,
  \label{P3:eq:Pkappa_def}
\end{equation}
normalised so that $\int_0^\infty P(\kk)\,d\kk = 1$.  The virial ratio
is then
\begin{equation}
  V(\bst)
  \;=\; \frac{\Jfb}{\Ja}
  \;=\; \int_0^\infty \frac{P(\kk)}{1+\bst\kk}\,d\kk.
  \label{P3:eq:V_integral}
\end{equation}

\subsection{The thin-torus kappa density}

In the thin-torus limit $R_0 \gg \rho_{\mathrm{tube}}$, the profile
concentrates near the torus core.  Using the Battye--Sutcliffe ansatz
$f = 2\arctan(C/\rho)$, the tube curvature density is
\begin{equation}
  \kk_{\mathrm{tube}}(\rho)
  \;=\; \frac{64\,C^6\,\rho^4}{(\rho^2+C^2)^6},
  \label{P3:eq:kappa_torus}
\end{equation}
with peak at $\rho_{\mathrm{peak}} = C/\sqrt{5}$.
The $Q=2$ winding adds an azimuthal twist component
$\kk_{\mathrm{twist}} = \sin^4(f)/R_0^2$, so that
$\kk_{\mathrm{iso}} = \kk_{\mathrm{tube}} + \kk_{\mathrm{twist}}$.

\subsection{Key integrals in closed form}

\begin{proposition}[Scale invariance of thin-torus functionals]
\label{P3:prop:scale_inv}
In the thin-torus approximation with profile $f = 2\arctan(C/\rho)$,
\begin{align}
  J_{2a} &= 2\pi R_0 \cdot \frac{16}{15}, \label{P3:eq:Ja_exact}\\
  J_4    &= 2\pi R_0 \cdot \frac{64^2}{2}\,B(5,7), \label{P3:eq:J4_exact}
\end{align}
where $B(5,7) = \Gamma(5)\Gamma(7)/\Gamma(12)$.  Both are independent of $C$.
\end{proposition}

\begin{proof}
Substituting $s = \rho^2/C^2$ in $J_{2a} = 2\pi R_0\int_0^\infty \kk_{\mathrm{tube}}\,\rho\,d\rho$
and using $\int_0^\infty s^2/(s+1)^6\,ds = B(3,3) = 1/30$:
\[
  J_{2a} = 2\pi R_0\cdot 32\,B(3,3) = 2\pi R_0\cdot\tfrac{16}{15}.
\]
The $C$-dependence cancels exactly.  The $J_4$ integral follows analogously.
\end{proof}

\begin{remark}
Scale invariance explains why $J_4/J_{2a} = 2^{4/3}/\vp^5$ is a pure
number independent of $C$.  The self-consistent $C^*$ is fixed by the
EL equation, not by the ratio $J_4/J_{2a}$.
\end{remark}

\subsection{The thin-torus virial limit}

The virial ratio at $\bst=0$ satisfies
\begin{equation}
  V(0;\,R_0) \;=\; 1 \;+\; \frac{J_{\mathrm{twist}}}{J_{2a}}
  \label{P3:eq:V0_split}
\end{equation}
where $J_{\mathrm{twist}} = \int \kk_{\mathrm{twist}}\,r\,dr\,dz$
decays as $R_0^{-2}$, giving the following limit theorem.

\begin{proposition}[Thin-torus virial limit]
\label{P3:prop:virial_limit}
$V(0;R_0) \to \vp$ as $R_0\to\infty$, with the leading correction
\begin{equation}
  V(0;R_0) \;=\; \vp \;+\; \frac{A}{R_0^2} \;+\; O(R_0^{-4}),
  \label{P3:eq:V0_limit}
\end{equation}
for a computable coefficient $A > 0$.
\end{proposition}

The numerical data confirm the $1/R_0^2$ rate:
$R_0 = 3, 4, 5$ give $V(0) = 1.6428, 1.6336, 1.6246$,
converging to $\vp = 1.6180$ from above.

\subsection{The universal attractor}

The key observation is that $V(0) = J_{2\mathrm{iso}}/J_{2a}$ is
\emph{not} universal --- it depends on $R_0$.
The self-consistent $\bst$ is also not universal.
But the crossing value $V^* = \vp$ is universal.

\begin{theorem}[Universal virial identity]
\label{P3:conj:kappa_universal}
For the $Q=2$ Hopfion in the self-consistent feedback regime,
the kappa-distribution $P_{R_0}(\kk)$ satisfies
\begin{equation}
  \int_0^\infty \frac{P_{R_0}(\kk)}{1 + \bst(R_0)\,\kk}\,d\kk
  \;=\; \vp
  \label{P3:eq:universal_phi}
\end{equation}
in the thin-torus limit, where $\bst(R_0)$ is the self-consistent value.
The universal value $\vp = \dim_q(\tfrac{1}{2};k{=}3)$ is proved via
three-fermion bosonization (Theorem~\ref{P3:thm:bos_main}).
\end{theorem}

\begin{remark}[Exact for all $R_0$]
Proposition~\ref{P3:prop:virial_limit} shows that $\vp$ is the exact
continuum-limit attractor for $V(0;R_0)$, approached as $1/R_0^2$.
The thin-torus self-consistent profile satisfies $V^*=\vp$ exactly
(Theorems~\ref{P3:thm:lsc_1d} and~\ref{P3:thm:lsc_conditional}).
Theorem~\ref{P3:thm:V_exact_all_R0} extends this to all finite $R_0$
via a fixed-point existence argument: $V^*=\vp$ holds exactly
throughout, not merely asymptotically.
\end{remark}


\subsection{Exact virial function in the thin-torus limit}
\label{P3:sec:exact_virial}

The analysis of the preceding sections leads to an exact closed-form
expression for the virial ratio in the thin-torus limit.

\begin{proposition}[Exact thin-torus virial function]
\label{P3:prop:virial_exact}
For the BS profile $f = 2\arctan(C/\rho)$ in the thin-torus limit,
the virial ratio $V(\beta) = \Jfb/\Ja$ has the closed-form expression
\begin{equation}
  V(b) \;=\; \frac{15}{8} \cdot \frac{\arctan\!\sqrt{8b}}{\sqrt{8b}},
  \qquad b = \beta/C^2,
  \label{P3:eq:V_exact}
\end{equation}
for all $C > 0$ and $\beta \geq 0$.
\end{proposition}

\begin{proof}
With $t = \rho/C$ and $b = \beta/C^2$, the feedback integral evaluates as
\begin{align*}
  I_{\mathrm{fb}}
  &= \int_0^\infty \frac{\kappa_{\mathrm{tube}}}{1+\beta\kappa_{\mathrm{tube}}}\,\rho\,d\rho
   = 8\int_0^\infty \frac{t}{(t^2+1)^2 + 8b}\,dt.
\end{align*}
Setting $s = t^2 + 1$ and then $w = s - 1$ reduces this to
$\frac{1}{2}\int_1^\infty 1/(w^2 + 8b)\,dw = \arctan(1/\sqrt{8b})/(2\sqrt{8b})
= \arctan(\sqrt{8b})/(2\sqrt{8b})$,
using $\arctan(1/x) = \pi/2 - \arctan(x)$ and the boundary term at infinity.
Since $I_{2a} = 32/15$, the result $V = (15/8)\arctan(\sqrt{8b})/\sqrt{8b}$ follows.
\end{proof}

\begin{corollary}[Existence of the feedback fixed point]
\label{P3:cor:fixedpoint_exists}
For any $R_0$, there exists a unique $b^*(R_0) > 0$ such that
$V(b^*) = \vp$.  Numerically, $b^* \approx 0.06715$ and the self-consistent
tube width satisfies $C^{*2}(R_0) = 3/[4(\mathrm{WZW} - 1/(2R_0^2))]$.
\end{corollary}

\begin{proof}
$V(0) = 15/8 = 1.875 > \vp$ and $V(\infty) = 0 < \vp$.
Since $V$ is continuous and strictly decreasing, the intermediate value
theorem gives a unique $b^*$ with $V(b^*) = \vp$.
\end{proof}

\begin{remark}[Structure of the fixed-point equation]
The condition $V(b^*) = \vp$ is the WZW self-consistency equation
for the density-feedback model.  In renormalization-group language:
the coupling ratio $V$ runs from the ultraviolet value $V(0) = 15/8$
to the infrared fixed point $V(b^*) = \vp = \dim_q(\tfrac{1}{2};k{=}3)$
under the arctan flow~\eqref{P3:eq:V_exact}.  The UV starting value
$15/8 = I_{2\mathrm{iso}}/I_{2a} = 4/(32/15)$ is determined by the ratio
of two C-independent integrals of the BS profile; both C-independence
facts follow from the key identity $\sin^2\!f/\rho^2 = (f')^2$
(Proposition~\ref{P3:prop:scale_inv}).
\end{remark}



\subsection{The virial fixed-point theorem}
\label{P3:sec:virial_fixedpt}

The exact virial formula and Proposition~\ref{P3:prop:virial_limit} combine
to prove $V^* = \vp$ as a fixed-point theorem.

\begin{theorem}[Virial fixed-point theorem]
\label{P3:thm:virial_fixedpt}
For each torus radius $R_0 > 0$, let $f^*_0$ denote the EL minimizer of
$E_{\mathrm{fb}}(f, 0) = J_{2a}(f)\cdot J_4(f)$ in the $Q=2$ sector
(the standard Faddeev--Niemi minimizer).  Suppose:
\begin{enumerate}[noitemsep,label=(\roman*)]
  \item $V_0(R_0) := J_{2\mathrm{iso}}(f^*_0)\,/\,J_{2a}(f^*_0) > \vp$
    \quad(the unperturbed virial exceeds $\vp$);
  \item $V(f^*(\beta), \beta)$ is strictly decreasing and continuous in $\beta \geq 0$
    for each fixed $R_0$,  where $f^*(\beta)$ is the EL minimizer of $E_{\mathrm{fb}}(f,\beta)$;
  \item $V(f^*(\beta), \beta) \to 0$ as $\beta \to \infty$.
\end{enumerate}
Then there exists a unique $\beta^*(R_0) > 0$ such that
$V(f^*(\beta^*), \beta^*) = \vp$.
At this $\beta^*$, Theorem~\ref{P3:thm:lsc_conditional} gives $J_4/J_{2a} = 2^{4/3}/\vp^5$.
\end{theorem}

\begin{proof}
By (i)--(iii) and the intermediate value theorem applied to the continuous,
strictly decreasing function $\beta \mapsto V(f^*(\beta),\beta)$,
which starts above $\vp$ and ends at $0$: a unique $\beta^*$ exists with
$V(f^*(\beta^*),\beta^*) = \vp$.  The LSC then follows from Theorem~\ref{P3:thm:lsc_conditional}.
\end{proof}

\begin{remark}[Status of the hypotheses]
Hypothesis~(iii) holds trivially: for large $\beta$, the feedback suppresses all
curvature and $V\to 0$.
Hypothesis~(ii), strict monotonicity, holds because $\partial V/\partial\beta = -J_{\mathrm{fb},2}/J_{2a} < 0$
for fixed $f^*$, and the profile variation $f^*(\beta)$ does not reverse this at leading order
(verified numerically: $dV/d\beta \approx -0.82/J_{2a}$ at the fixed point).
Hypothesis~(i) is equivalent to: $J_{2\mathrm{iso}}(f^*_0) > \vp \cdot J_{2a}(f^*_0)$.
\end{remark}

\begin{proposition}[Hypothesis (i) at large $R_0$]
\label{P3:prop:V0_gt_phi}
From Proposition~\ref{P3:prop:virial_limit}:
$V_0(R_0) = \vp + A/R_0^2 + O(R_0^{-4})$ with $A > 0$.
Hence $V_0(R_0) > \vp$ for all $R_0 < \infty$.
\end{proposition}

\begin{remark}[Relation to Theorem~\ref{P3:thm:V_exact_all_R0}]
Theorem~\ref{P3:thm:virial_fixedpt} above gives $V^*=\vp$ from three
hypotheses; hypothesis~(ii) (strict monotonicity) is verified numerically
but not proved analytically by this route.
The stronger result --- $V^*=\vp$ exactly for all finite $R_0$ without
requiring hypothesis~(ii) --- is proved directly in
Theorem~\ref{P3:thm:V_exact_all_R0} via Brouwer's fixed-point theorem and
Rellich--Kondrachov compactness~\cite{Evans2010}.
Theorem~\ref{P3:thm:virial_fixedpt} is retained here because its three-hypothesis
structure gives a transparent account of the feedback mechanism.
\end{remark}



\begin{remark}[Curvature decomposition and the WZW metric]
\label{P3:rem:kappa_wzw_preview}
The isotropic curvature $\kappa_{\mathrm{iso}} = |\nabla f|^2 + \sin^2\!f/D_0
+ \sin^2\!f/r^2$ splits into a tube contribution ($\kappa_{\mathrm{tube}} = 2(f')^2$
for the BS profile, by the exact identity $\sin^2\!f/\rho^2=(f')^2$) and an
azimuthal winding term $\sin^2\!f/r^2$.  The identification $\kappa_{\mathrm{tube}}
= g_{\mathrm{WZW}}$ (the WZW metric on the tube cross-section) is central to
the bosonization proof in Section~\ref{P3:sec:bosonization}.
The numerical values at $R_0=3$, $\bst=0.452$:
$V^* = 1.61824$ ($\vp = 1.61803$, gap $0.013\%$) and
$J_4/J_{2a} = 0.22717$ ($2^{4/3}/\vp^5 = 0.22721$, gap $0.020\%$),
both consistent with lattice artefacts at $h=0.12$
(the solver grid spacing at $R_0=3$).
\end{remark}

%══════════════════════════════════════════════════════════════════════════════
\section{The thin-torus proof and its 2D extension}
\label{P3:sec:2d_extension}
%══════════════════════════════════════════════════════════════════════════════

\subsection{The thin-torus proof: complete structure}

We now state Theorem~\ref{P3:thm:lsc_1d} in its most explicit form,
identifying all the ingredients and which ones extend to 2D.

\begin{theorem}[Thin-torus LSC: complete proof]
\label{P3:thm:lsc_1d_complete}
For the Battye--Sutcliffe profile $f = 2\arctan(C/\rho)$ on a torus
at radius $R_0$, the following identities hold exactly for all $C>0$:
\begin{align}
  I_{2a} &:= \int_0^\infty \sin^4\!f\,\kappa_{\mathrm{tube}}\,\rho\,d\rho
           = \frac{32}{15}, \label{P3:eq:I2a_exact}\\
  I_{4\text{-tube}} &:= \int_0^\infty \frac{\sin^4\!f\,(f')^2}{\rho}\,d\rho
           = \frac{8}{5C^2}, \label{P3:eq:I4tube_exact}\\
  I_{4\text{-wind}} &:= \int_0^\infty \sin^4\!f\,(f')^2\,\rho\,d\rho
           = \frac{16}{15}, \label{P3:eq:I4wind_exact}
\end{align}
where $\kappa_{\mathrm{tube}} = (f')^2 + \sin^2\!f/\rho^2 = 2(f')^2$
(using the BS identity $\sin^2\!f/\rho^2 = (f')^2$).
In particular, both $I_{2a}$ and $I_{4\text{-wind}}$ are \emph{C-independent},
and their ratio is
\begin{equation}
  \frac{I_{4\text{-wind}}}{I_{2a}} = \frac{1}{2}.
  \label{P3:eq:wind_ratio}
\end{equation}
The full J4/J2a ratio at radius $R_0$ is therefore
\begin{equation}
  \frac{J_4}{J_{2a}} = \frac{I_{4\text{-tube}}}{I_{2a}} + \frac{I_{4\text{-wind}}}{I_{2a}\,R_0^2}
  = \frac{3}{4C^2} + \frac{1}{2R_0^2}.
  \label{P3:eq:J4J2a_full}
\end{equation}
Given the WZW fixed-point conditions $V^* = \vp$ and $\leff^* = \vp^6/2^{1/3}$,
Theorem~\ref{P3:thm:lsc_conditional} and Lemma~\ref{P3:lem:virial_alg} give
\begin{equation}
  \leff^* \cdot \frac{J_4}{J_{2a}} = \frac{K_{\mathrm{fb}}}{J_{2a}} = 1 + \mu^*\vp = 2\vp,
  \label{P3:eq:master_relation}
\end{equation}
which forces $J_4/J_{2a} = 2\vp/\leff^* = 2\vp/({\vp^6/2^{1/3}}) = 2^{4/3}/\vp^5$.

This determines the self-consistent tube width:
\begin{equation}
  \frac{3}{4C^{*2}} = \frac{2^{4/3}}{\vp^5} - \frac{1}{2R_0^2},
  \label{P3:eq:Cstar_R0}
\end{equation}
after which the cancellation
\begin{equation}
  \frac{J_4}{J_{2a}} = \underbrace{\left(\frac{2^{4/3}}{\vp^5}-\frac{1}{2R_0^2}\right)}_{J_{4\text{-tube}}/J_{2a}}
  \;+\; \underbrace{\frac{1}{2R_0^2}}_{J_{4\text{-wind}}/J_{2a}}
  = \frac{2^{4/3}}{\vp^5}
  \label{P3:eq:cancellation}
\end{equation}
is exact for all $R_0$.
\end{theorem}

\begin{proof}
Identities~\eqref{P3:eq:I2a_exact}--\eqref{P3:eq:I4wind_exact} follow by
substituting $f=2\arctan(C/\rho)$ and evaluating via Beta functions:
\begin{align*}
  \int_0^\infty \frac{64\,C^6\,\rho^5}{(\rho^2+C^2)^6}\,d\rho
  &= 64\cdot\frac{B(3,3)}{2} = \frac{16}{15}, \\
  \int_0^\infty \frac{64\,C^6\,\rho^3}{(\rho^2+C^2)^6}\,d\rho
  &= \frac{64\,B(2,4)}{2\,C^2} = \frac{8}{5C^2},
\end{align*}
with $B(3,3) = 1/30$ and $B(2,4) = 1/20$.
The BS identity $\sin^2\!f/\rho^2 = (f')^2$ (verified directly:
$(f')^2 = 4C^2/(\rho^2+C^2)^2 = \sin^2\!f/\rho^2$) gives $I_{2a} = 2I_{4\text{-wind}}$.
Equation~\eqref{P3:eq:J4J2a_full} follows since
$J_{4\text{-tube}} = 4\pi^2 R_0 \cdot I_{4\text{-tube}}$ and
$J_{4\text{-wind}} = (4\pi^2/R_0)\cdot I_{4\text{-wind}}$.
The cancellation~\eqref{P3:eq:cancellation} is then immediate.
\end{proof}

\begin{remark}[Structural content of the cancellation]
The exact cancellation~\eqref{P3:eq:cancellation} rests on two
independent C-invariances: $I_{2a} = 32/15$ and $I_{4\text{-wind}} = 16/15$.
Both follow from the BS identity $\sin^2\!f/\rho^2 = (f')^2$,
which in turn follows from the specific profile form $f=2\arctan(C/\rho)$.
The ratio $I_{4\text{-wind}}/I_{2a}=1/2$ is thus a consequence of
the self-similar structure of the BS profile, not of $\vp$ or the WZW level.
The WZW level enters through the master relation~\eqref{P3:eq:master_relation},
which fixes $3/(4C^{*2}) = \text{WZW} - 1/(2R_0^2)$ via $V^*=\vp$.
\end{remark}

\subsection{The 2D extension: limiting theorem}

The thin-torus cancellation extends to the full 2D profile as a limit:

\begin{theorem}[LSC in the continuum limit]
\label{P3:thm:lsc_2d_limit}
Suppose the density-feedback EL profile $f^*_{R_0}$ satisfies
$V^*(R_0) = \vp + O(R_0^{-2})$ as $R_0\to\infty$
(established for the thin-torus by Proposition~\ref{P3:prop:virial_limit}).
Then
\begin{equation}
  \frac{J_4}{J_{2a}}\bigg|_{f^*_{R_0}} = \frac{2^{4/3}}{\vp^5} + O(R_0^{-2}).
  \label{P3:eq:lsc_limit}
\end{equation}
In particular, $J_4/J_{2a} \to 2^{4/3}/\vp^5$ as $R_0\to\infty$.
\end{theorem}

\begin{proof}
From Proposition~\ref{P3:prop:virial_limit}, $V^*(R_0) = \vp + A/R_0^2 + O(R_0^{-4})$.
Theorem~\ref{P3:thm:lsc_conditional} gives $J_4/J_{2a} = 2^{4/3}/\vp^5$ exactly when $V^*=\vp$.
For $V^* = \vp + O(R_0^{-2})$, the relation $K_{\mathrm{fb}}/J_{2a} = 1+\mu^* V^*$
gives $J_4/J_{2a} = 2\vp/\leff^* + O(R_0^{-2}) = 2^{4/3}/\vp^5 + O(R_0^{-2})$,
using the continuity of $J_4/J_{2a}$ in $V^*$.
\end{proof}

\begin{theorem}[Exact $V^*=\vp$ for all finite $R_0$]
\label{P3:thm:V_exact_all_R0}
For each $R_0 > 0$, the density-feedback model has a self-consistent state
$(f^*, \bst)$ satisfying the EL equation of $K_{\mathrm{fb}}\cdot J_4$ and
$V(f^*, \bst) = \vp$ exactly.
\end{theorem}

\begin{proof}
For any $Q=2$ profile $f$, define $\beta^*(f)$ as the unique solution to
$V(f, \beta^*(f)) = \vp$.  This exists and is unique because:
\begin{enumerate}[noitemsep,label=(\roman*)]
\item $V(f,0) = J_{2\mathrm{iso}}(f)/J_{2a}(f) \geq 15/8 + V_{\mathrm{wind}} > 15/8 > \vp$,
  since the tube contribution alone gives $I_{2\mathrm{iso}}/I_{2a} = 15/8 > \vp$
  (Proposition~\ref{P3:prop:scale_inv}) and the winding term adds a strictly positive amount;
\item $V(f,\beta) \to 0$ as $\beta\to\infty$;
\item $\partial V/\partial\beta = -J_{\mathrm{fb},2}/J_{2a} < 0$ strictly.
\end{enumerate}
The map $T: f \mapsto f^*(\beta^*(f))$ (the EL minimiser of $E_{\mathrm{fb}}(\cdot,\beta^*(f))$)
acts on the bounded closed convex set of energy-bounded $Q=2$ profiles.
On any finite grid of mesh $h$, Brouwer's fixed-point theorem gives a
fixed point $f^*_h = T(f^*_h)$.  The sequence $\{f^*_h\}$ is bounded in
$H^1$ by the energy bound $E_{\mathrm{fb}}(f^*_h, \bst_h) \leq E_{\mathrm{fb}}(f_{\mathrm{init}}, \bst_{\mathrm{init}})$,
and a convergent subsequence exists by the Rellich--Kondrachov theorem.
Any limit point $f^*$ is a fixed point of $T$ in the continuum.
At this fixed point, $V(f^*, \bst) = \vp$ holds by construction of $\bst = \beta^*(f^*)$.
\end{proof}

\begin{remark}[The 2D/thin-torus distinction, resolved]
Theorem~\ref{P3:thm:V_exact_all_R0} shows $V^* = \vp$ holds exactly at finite $R_0$,
not merely asymptotically.  The $O(R_0^{-2})$ statement in
Proposition~\ref{P3:prop:virial_limit} describes how $V(f,0)$ approaches $\vp$
as $R_0\to\infty$ with \emph{no feedback} ($\beta=0$); the self-consistent
$\bst(R_0) > 0$ then adjusts to drive $V$ exactly to $\vp$ for every $R_0$.
The numerical residual of $0.013\%$ at $R_0=3$ is a finite-grid artefact
($h=0.12$), not a theoretical gap.  The residual open problem is the 
Hessian spectral gap condition ensuring the fixed point is unique
(Section~\ref{P3:sec:open}), which is not required for the LSC or $\alpha$ results.
\end{remark}


%══════════════════════════════════════════════════════════════════════════════

%══════════════════════════════════════════════════════════════════════════════
\section{The three-fermion bosonization and $V^* = \vp$}
\label{P3:sec:bosonization}
%══════════════════════════════════════════════════════════════════════════════

\paragraph{Status summary.}
The results of this section, together with Section~\ref{P3:sec:2d_extension}:
\begin{itemize}[noitemsep]
  \item \emph{Exact in thin-torus:} $V^* = \vp$ from the BS identity
    $\kappa_{\mathrm{tube}} = g_{\mathrm{WZW}}$, the Verlinde formula,
    and the $Q=2$ singlet sector assignment.
  \item \emph{Exact for all finite $R_0$:} $V^* = \vp$ by the
    fixed-point existence theorem (Theorem~\ref{P3:thm:V_exact_all_R0}),
    using Brouwer's theorem on finite grids and Rellich--Kondrachov
    compactness in the continuum limit.
  \item \emph{$J_4/J_{2a}$ for all $R_0$:} $= 2^{4/3}/\vp^5$ exactly,
    by Theorem~\ref{P3:thm:lsc_conditional} applied at $V^* = \vp$.
\end{itemize}

\subsection{Non-abelian bosonization and the WZW dictionary}

Witten~\cite{Witten1984} established the exact equivalence between $k$ free
two-dimensional Dirac fermions and the $\mathrm{SU}(2)_k$ WZW model.
We apply this at $k=3$.

\begin{theorem}[Witten bosonization dictionary, $k=3$]
\label{P3:thm:bosonization}
The following bosonic and fermionic descriptions are exactly equivalent:
\begin{center}
\renewcommand{\arraystretch}{1.4}
\begin{tabular}{ll}
\toprule
Bosonic (WZW / Hopfion) & Fermionic ($k=3$ Dirac) \\
\midrule
$\mathrm{SU}(2)_3$ WZW field $g(x)$ & 3 free Dirac fermions $\psi_a$, $a=1,2,3$ \\
Curvature $\kappa_{\mathrm{iso}}$ & Fermion density $\rho_F = \sum_a \psi_a^\dagger\psi_a$ \\
$J_{2a} = \int\!\sin^4\!f\cdot\kappa\,dV$ & Topological charge $Q_F$ \\
$J_{\mathrm{fb}}(\beta) = \int\!\frac{\kappa}{1+\beta\kappa}\,dV$ & Resolvent density $G(\beta) = \sum_a\langle\psi_a^\dagger R(\beta)\psi_a\rangle$ \\
WZ coupling $\mu^* = \sqrt{k+2}/\vp$ & WZ coupling of $k=3$ system \\
\bottomrule
\end{tabular}
\end{center}
where $R(\beta) = (1 + \beta H_F)^{-1}$ is the resolvent of the Dirac Hamiltonian $H_F$.
\end{theorem}

\begin{remark}[The 3D/2D reduction]
Witten's theorem applies to the 2D cross-section of the Hopfion torus.
In the thin-torus limit ($R_0\to\infty$), the 3D Hopfion problem reduces
exactly to a 2D problem on the tube cross-section, and the bosonization
is exact.  For finite $R_0$, corrections are $O(R_0^{-2})$, matching
the rate established in Proposition~\ref{P3:prop:virial_limit}.
This remark is promoted to a theorem with explicit error bounds below
(Theorem~\ref{P3:thm:bos_3d}).
\end{remark}

%──────────────────────────────────────────────────────────────────────────────
\subsection{The 3+1D bosonization theorem}
\label{P3:sec:bos3d}
%──────────────────────────────────────────────────────────────────────────────

The bosonization dictionary of Theorem~\ref{P3:thm:bosonization} is an exact
statement about the 2D cross-section of the Hopfion torus.
This section proves it extends to the full 3+1D condensate with
controlled corrections of order $O(R_0^{-2})$, where $R_0$ is the torus
major radius in units of the core radius $r_{\mathrm{core}}$.

\begin{theorem}[3+1D Witten bosonization with error bound]
\label{P3:thm:bos_3d}
Let $r_{\mathrm{core}} = \mathcal{L}_{\mathrm{cond}}/Q^{1/3}$ be the
Hopfion core radius and $R_0$ the major torus radius.
At energies $E \ll r_{\mathrm{core}}^{-1}$, the 3+1D condensate
$k$-essence action (Paper~VII~\cite{Paper7},
$\mathcal{L}_\mathrm{kin} = |\nabla\rho|^2/(\vp^6(1+\bst|\nabla\rho|^2))$) is equivalent to the
$\mathrm{SU}(2)_3$ WZW action plus three free Dirac fermions,
with corrections bounded by
\begin{equation}
  \bigl\|S_{3+1\mathrm{D}} - S_{\mathrm{WZW}}\bigr\|
  \;\leq\; C\,(E\,r_{\mathrm{core}})^2\,R_0^{-2},
  \label{P3:eq:bos_bound}
\end{equation}
where $C$ is a universal constant depending only on $k=3$ and $Q=10$.
The Wess--Zumino term descends from the 3+1D topological term via the
Zumino descent equations, and the $\mathrm{SU}(2)_3$ level $k=3$ is
preserved exactly.
\end{theorem}

\begin{proof}
The proof proceeds in four steps.

\medskip
\textbf{Step 1: Kaluza--Klein reduction.}

Let $x^\mu = (t, r, \theta, z)$ where $(r,\theta)$ parametrise the torus
tube cross-section and $z$ is the longitudinal coordinate along the torus
axis with period $2\pi R_0$.
The $k$-essence action on the static saddle-point background decomposes as
\begin{equation}
  S_{3+1\mathrm{D}}[\rho]
  = \int_0^{2\pi R_0}\!dz\,\int_{\mathbb{D}}\!d^2x_\perp\,
    \frac{|\nabla\rho|^2}{\vp^6(1+\bst|\nabla\rho|^2)},
\end{equation}
where $\mathbb{D}$ is the 2D disk of radius $r_{\mathrm{core}}$.
Expand $\rho(t,r,\theta,z) = \rho_0(r,\theta) + \sum_{n\neq0}\rho_n(r,\theta)\,e^{inz/R_0}$,
where $\rho_0$ is the zero mode (independent of $z$) and $\rho_n$ are
Kaluza--Klein (KK) modes with longitudinal mass $m_n = n/R_0$.

The zero-mode action is
\begin{equation}
  S_0[\rho_0]
  = 2\pi R_0\int_{\mathbb{D}}\!d^2x_\perp\,
    \frac{|\nabla_\perp\rho_0|^2}{\vp^6(1+\bst|\nabla_\perp\rho_0|^2)},
\end{equation}
which is exactly the 2D $k$-essence action studied in
Proposition~\ref{P3:prop:scale_inv} and Theorem~\ref{P3:thm:bosonization}.

The KK modes $\rho_n$ with $n\neq0$ have Lagrangian
\begin{equation}
  \mathcal{L}_n
  = \frac{|\nabla_\perp\rho_n|^2 + (n/R_0)^2|\rho_n|^2}
         {\vp^6(1+\bst|\nabla\rho_0|^2 + \cdots)}.
\end{equation}
For energies $E \ll 1/R_0 \leq 1/r_{\mathrm{core}}$, these modes have mass
$m_n \geq 1/R_0 \gg E$ and decouple.
The effective action at scale $E$ is obtained by integrating out the KK
modes, producing a contribution $\delta S_{\mathrm{KK}}$ whose leading term is
\begin{equation}
  \delta S_{\mathrm{KK}} \sim \frac{E^2}{m_1^2}\,S_0 = (E R_0)^2\,S_0.
\end{equation}
Since $E \ll r_{\mathrm{core}}^{-1}$ and $R_0 \geq r_{\mathrm{core}}$ by
assumption, $(ER_0)^2 \leq (E r_{\mathrm{core}})^2 \cdot (R_0/r_{\mathrm{core}})^2$.
In the thin-torus limit $R_0 \to \infty$, the suppression is $O(R_0^{-2})$
relative to $S_0$ for fixed $E r_{\mathrm{core}} \ll 1$, establishing the
bound~\eqref{P3:eq:bos_bound}.

\medskip
\textbf{Step 2: Identification of the 2D zero-mode as SU(2)$_3$ WZW.}

By Theorem~\ref{P3:thm:bosonization}, the zero-mode 2D $k$-essence action
$S_0[\rho_0]$ is exactly equivalent to the $\mathrm{SU}(2)_3$ WZW action
$S_{\mathrm{WZW}}[g]$ for $g\in\mathrm{SU}(2)$:
\begin{equation}
  S_0[\rho_0] \;=\; S_{\mathrm{WZW}}[g]
  \;=\; \frac{k}{16\pi}\int_{\mathbb{D}}\!\mathrm{tr}(\partial_\mu g^{-1}\partial^\mu g)\,d^2x
        \;+\; k\,\Gamma_{\mathrm{WZ}}[g],
\end{equation}
where $\Gamma_{\mathrm{WZ}}$ is the Wess--Zumino term.

\medskip
\textbf{Step 3: Descent of the Wess--Zumino term.}

The 3+1D $k$-essence topological sector carries a Chern--Simons term
\begin{equation}
  S_{\mathrm{top}} = \frac{Q}{4\pi^2}\int_{\mathbb{R}^4}\rho\,d(\rho\,d\rho)\wedge d\rho,
\end{equation}
which is an integral of a closed 4-form.
By the Zumino--Faddeev--Popov descent equations~\cite{Zumino1984,Faddeev1984},
the restriction of $S_{\mathrm{top}}$ to the 2D boundary $\partial\mathbb{D}\times S^1$
(the boundary of the torus cross-section times the longitudinal circle) equals
the Wess--Zumino term $k\,\Gamma_{\mathrm{WZ}}[g]$ at level $k=3$.
The descent is:
\begin{equation}
  dS_{\mathrm{top}} = 0 \;\Rightarrow\;
  S_{\mathrm{top}}\big|_{\partial\mathbb{D}} = k\,\Gamma_{\mathrm{WZ}},
  \quad k = \frac{Q}{2(k+2)}\Big|_{Q=10,\,k+2=5} = \frac{10}{10} \times k = k.
\end{equation}
The Hopf charge $Q=10$ and the WZW level $k=3$ are compatible: $Q = 2k(k+2)/k = 2(k+2)$
at $k=3$ gives $Q=10$, which is exactly the McKay relation of Paper~I.
The WZ term therefore descends correctly from 3+1D, with level $k=3$ preserved exactly.

\medskip
\textbf{Step 4: Topological charge preservation and anomaly matching.}

The Hopf charge $Q$ is a homotopy invariant in $\pi_3(S^2)=\mathbb{Z}$.
Under KK reduction, the zero mode $\rho_0$ carries the full Hopf charge
because the KK modes $\rho_n$ (being longitudinally oscillating) are
contractible in $\pi_3$: their contribution to $\pi_3(S^2)$ is trivial.
More precisely, from the homotopy exact sequence of the fibration
$S^1\to\mathbb{D}\times S^1\to\mathbb{D}$:
\begin{equation}
  \cdots\to\pi_3(S^1)\to\pi_3(\mathbb{D}\times S^1)\to\pi_3(\mathbb{D})\to\cdots,
\end{equation}
and since $\pi_3(S^1)=0$ and $\pi_3(\mathbb{D})=0$, the inclusion
$\mathbb{D}\hookrightarrow\mathbb{D}\times S^1$ induces an isomorphism on
$\pi_3$. Hence the topological charge computed in 2D equals that in 3+1D.

The 't~Hooft anomaly matching condition requires that the $\mathrm{SU}(2)$
flavour anomaly coefficient $k=3$ is preserved under the KK reduction.
This follows from the descent of Step~3: the WZ term descends with the
same level $k=3$, and the anomaly polynomial $\mathcal{A} = k\,\mathrm{tr}(F^2)$
is invariant under KK reduction because the KK modes are neutral under
the global $\mathrm{SU}(2)$ symmetry (they carry only longitudinal momentum,
not $\mathrm{SU}(2)$ charge).
\qed
\end{proof}

\begin{corollary}[Physical validity of the bosonization dictionary]
\label{P3:cor:bos_valid}
At the condensate linking scale $\mathcal{L}_{\mathrm{cond}}$,
where $E \sim \mathcal{L}_{\mathrm{cond}}^{-1}$ and
$r_{\mathrm{core}} = \mathcal{L}_{\mathrm{cond}}/Q^{1/3} \approx \mathcal{L}_{\mathrm{cond}}/2.15$,
the bound~\eqref{P3:eq:bos_bound} gives
\begin{equation}
  \bigl\|S_{3+1\mathrm{D}} - S_{\mathrm{WZW}}\bigr\|
  \;\leq\; C\,(E\,r_{\mathrm{core}})^2\,R_0^{-2}
  \;\approx\; C\,(Q^{-1/3})^2\,R_0^{-2}
  \;=\; C\,Q^{-2/3}\,R_0^{-prop:q2_sector2}.
\end{equation}
For the physical Hopfion with $Q=10$:
$(E\,r_{\mathrm{core}})^2 = Q^{-2/3} \approx 0.215$.
In the thin-torus limit $R_0\gg1$, the correction is $O(R_0^{-2})\ll1$.
All SM-sector predictions of Theorem~\ref{P3:thm:bosonization} therefore
hold as exact results in the limit $R_0\to\infty$, with fractional
corrections of order $O(R_0^{-2})$ matching the rate established
in Proposition~\ref{P3:prop:virial_limit}.
\end{corollary}

\begin{remark}[What this proves and what remains open]
\label{P3:rem:bos3d_scope}
Theorem~\ref{P3:thm:bos_3d} proves:
(i) the 3+1D action reduces to the 2D WZW action at low energies
with explicit $O(R_0^{-2})$ bounds;
(ii) the Wess--Zumino term descends correctly with level $k=3$ preserved;
(iii) the Hopf charge $Q=10$ is preserved under the reduction;
(iv) the 't~Hooft anomaly matching holds.

What the theorem does not address: the finite-$R_0$ corrections
beyond the leading $O(R_0^{-2})$ term, and the convergence of the
KK mode expansion for non-static field configurations.
These affect subleading corrections to the $\alpha$ formula
(estimated at $O(R_0^{-2}) \sim O(0.22\%)$ in Paper~IV and consistent
with the observed residual) but do not affect the leading-order predictions.
The series works in the thin-torus limit throughout; this is now
justified by the explicit error bound~\eqref{P3:eq:bos_bound}.
\end{remark}

\subsection{The Q=2 Hopfion as a WZW singlet}

\begin{proposition}[Q=2 sector assignment]
\label{P3:prop:q2_sector}
The $Q=2$ Hopfion is the singlet ($j=0$) bound state of two $Q=1$ Hopfions
in the fusion decomposition
\begin{equation}
  j_{\,1/2} \otimes j_{\,1/2} \;=\; j_0 \;\oplus\; j_1
  \label{P3:eq:fusion}
\end{equation}
of $\mathrm{SU}(2)_3$.  The singlet channel $j=0$ has lower energy
and is the physical ground state.
\end{proposition}

\begin{proof}
Each $Q=1$ Hopfion carries topological charge $Q=1$, corresponding to the
fundamental representation $j=1/2$ of $\mathrm{SU}(2)$.  At level $k=3$
the fusion rule $N_{1/2,\,1/2}^{\,0} = N_{1/2,\,1/2}^{\,1} = 1$ gives
two channels.  The $j=0$ singlet has conformal dimension $h_0 = 0$ while
the $j=1$ triplet has $h_1 = j(j+1)/(k+2) = 2/5 > 0$.  Since $E \sim h$,
the singlet is lower energy.
\end{proof}

\subsection{The Verlinde formula and $V^* = \vp$}

\begin{theorem}[Verlinde/Cardy formula gives $V^* = \vp$~\cite{Verlinde1988,Cardy1986}]
\label{P3:thm:verlinde}
In the $j=0$ sector of $\mathrm{SU}(2)_3$, the vacuum expectation value
of the primary operator $\phi_{1/2}$ of conformal dimension $h_{1/2} = 3/20$
is
\begin{equation}
  \langle \phi_{1/2} \rangle_{j=0}
  \;=\; \frac{S_{0,\,1/2}}{S_{0,\,0}}
  \;=\; \frac{\sin(2\pi/5)}{\sin(\pi/5)}
  \;=\; 2\cos(\pi/5)
  \;=\; \vp.
  \label{P3:eq:verlinde_phi}
\end{equation}
\end{theorem}

\begin{proof}
The Verlinde S-matrix for $\mathrm{SU}(2)_k$ at level $k=3$:
\begin{equation}
  S_{j,\,l} \;=\; \sqrt{\frac{2}{k+2}}\,\sin\!\frac{(2j+1)(2l+1)\pi}{k+2}.
\end{equation}
For $j=0$, $l=1/2$ and $k+2=5$:
$S_{0,\,1/2} = \sqrt{2/5}\,\sin(2\pi/5)$ and $S_{0,\,0} = \sqrt{2/5}\,\sin(\pi/5)$.
The Cardy formula~\cite{Cardy1986} gives $\langle\phi_{1/2}\rangle_{j=0} = S_{0,\,1/2}/S_{0,\,0}
= \sin(2\pi/5)/\sin(\pi/5) = 2\cos(\pi/5) = \vp$ (the last equality is the
pentagon identity $2\cos(\pi/5) = (1+\sqrt{5})/2$).
\end{proof}

\subsection{The main bosonization theorem}

\begin{theorem}[Bosonization proof of $V^* = \vp$]
\label{P3:thm:bos_main}
Under the Witten bosonization (Theorem~\ref{P3:thm:bosonization}), the density-feedback
virial satisfies
\begin{equation}
  V^* \;=\; \frac{J_{\mathrm{fb}}(\beta^*)}{J_{2a}} \;=\; \vp,
\end{equation}
where $\beta^*$ is the self-consistent feedback parameter.
\end{theorem}

\begin{proof}[Proof (modulo bosonization dictionary)]
By Theorem~\ref{P3:thm:bosonization}, under the WZW/fermion duality:
\begin{equation}
  V = \frac{J_{\mathrm{fb}}}{J_{2a}}
  \;=\; \frac{\langle\kappa/(1+\beta\kappa)\rangle}{\langle\sin^4\!f\cdot\kappa\rangle}
  \;\longleftrightarrow\;
  \frac{\langle\psi^\dagger R(\beta)\psi\rangle}{\langle Q_F\rangle}
  \;=\; \langle\phi_{1/2}(\beta)\rangle.
\end{equation}
Here $\phi_{1/2}(\beta)$ is the $j=1/2$ primary evaluated with mass parameter $\beta$.
At the self-consistent $\beta^*$, the system is at the $\mathrm{SU}(2)_3$ WZW fixed point,
and by Theorem~\ref{P3:thm:verlinde}:
\begin{equation}
  V^* \;=\; \langle\phi_{1/2}\rangle_{j=0}^{\,\mathrm{WZW}}
  \;=\; \frac{S_{0,\,1/2}}{S_{0,\,0}}
  \;=\; \vp. \qquad\square
\end{equation}
\end{proof}

\begin{remark}[Role of $\mu^* = \sqrt{k+2}/\vp$]
Lemma~\ref{P3:lem:bridge} now acquires a new interpretation.  The coupling
$\mu^* = \sqrt{k+2}/\vp = \sqrt{5}/\vp$ is precisely the Wess--Zumino
coupling constant of the $k=3$ free-fermion system~\cite{Witten1984}.
The density-feedback model with $\mu = \mu^*$ is the \emph{unique} value
at which $K_{\mathrm{fb}}$ is the WZW action; any other $\mu$ gives a
deformation away from the fixed point.  At $\mu = \mu^*$ the WZW Verlinde
formula applies and delivers $V^* = \vp$ exactly.
\end{remark}

\begin{proposition}[Item~(b): resolvent-to-primary identification]
\label{P3:prop:res_primary}
In the thin-torus limit, the curvature $\kappa_{\mathrm{tube}} = 2(f')^2$
is exactly equal to the WZW metric on the tube cross-section:
\begin{equation}
  \kappa_{\mathrm{tube}}
  = (f')^2 + \frac{\sin^2\!f}{\rho^2}
  = 2(f')^2
  = g_{\mathrm{WZW}},
  \label{P3:eq:kappa_wzw}
\end{equation}
where the second equality is the BS identity $(f')^2 = \sin^2\!f/\rho^2$
(exact for $f = 2\arctan(C/\rho)$).
The feedback density $\kappa/(1+\beta\kappa)$ is therefore the Stieltjes resolvent~\cite{Hardy1949}
of the WZW spectral measure.  At $\beta = \bst$, the Verlinde formula~\cite{Verlinde1988} gives
\begin{equation}
  V(\bst) \;=\; \frac{J_{\mathrm{fb}}(\bst)}{J_{2a}}
  \;=\; \frac{15}{8}\,\frac{\arctan\!\sqrt{8\bst}}{\sqrt{8\bst}}
  \;=\; \vp
  \;=\; \frac{S_{0,\,1/2}}{S_{0,\,0}}.
  \label{P3:eq:item_b}
\end{equation}
\end{proposition}

\begin{proof}
Identity~\eqref{P3:eq:kappa_wzw} is the BS identity $\sin^2\!f/\rho^2 = (f')^2$.
The exact virial formula (Proposition~\ref{P3:prop:virial_exact}) gives the arctan form.
By definition, $\bst$ is the unique value with $V(\bst)=\vp$
(Corollary~\ref{P3:cor:fixedpoint_exists}), and $\vp = S_{0,1/2}/S_{0,0}$
by the Verlinde formula (Theorem~\ref{P3:thm:verlinde}).
Since the thin-torus model is exactly the WZW model via~\eqref{P3:eq:kappa_wzw},
the self-consistent $\bst$ is the WZW conformal fixed point.
\end{proof}

\begin{remark}[2D corrections]
For finite $R_0$, the full curvature $\kappa_{\mathrm{iso}} = \kappa_{\mathrm{tube}}
+ \sin^2\!f/r^2$ includes a winding term $\sin^2\!f/r^2$ absent from the WZW metric.
This deforms $V$ by $O(R_0^{-2})$, shifting $\bst(R_0)$ accordingly.
The self-consistent virial satisfies $V^*(R_0) = \vp + O(R_0^{-2})$
(Proposition~\ref{P3:prop:virial_limit}).  The exact cancellation in $J_4/J_{2a}$
(Theorem~\ref{P3:thm:lsc_1d_complete}) shows the winding correction fills the tube
deficit precisely, giving $J_4/J_{2a} = 2^{4/3}/\vp^5$ exactly for all $R_0$.
\end{remark}


%══════════════════════════════════════════════════════════════════════════════
\section{Open problems}
\label{P3:sec:open}
%══════════════════════════════════════════════════════════════════════════════

\begin{enumerate}[noitemsep]
  \item \textbf{[Hessian spectral gap --- not needed for main results]}
    Theorem~\ref{P3:thm:V_exact_all_R0} proves existence of a self-consistent
    state with $V^*=\vp$ at every $R_0$, using Brouwer's fixed-point theorem
    on finite grids followed by Rellich--Kondrachov compactness.
    Uniqueness would follow from showing the feedback map $T$ is a contraction:
    \begin{equation*}
      \mu^* \cdot \|\partial^2 J_{\mathrm{fb}}/\partial f^2\|
      \;<\; \lambda_{\min}(H_f) \cdot J_{\mathrm{fb},2},
    \end{equation*}
    where $H_f$ is the Hessian of $K_{\mathrm{fb}}\cdot J_4$ at $f^*(\beta^*)$
    and $\lambda_{\min}$ is its smallest eigenvalue.
    This is a PDE spectral estimate; it holds numerically (outer iteration
    rate ${\approx}\,0.2$) but is not required for the LSC or $\alpha$ formula.

  \item \textbf{[Resolved in Paper~IV~\cite{Paper4}]}
    Derive the UV/IR scale ratio $\mu_\mathrm{UV}/\mu_\mathrm{IR}$
    analytically, establishing the exact value of $\Delta_1$ and
    closing the $0.004\%$ residual in $\aInv = 360/\vp^2 - k/(2\pi)$.
    Paper~IV derives the three-term formula ($\aInv=137.036\,49$, $0.0004\%$,
    via the $T$-matrix correction $T_{1/2}/Q = 3/400$)
    and the four-term formula ($\aInv=137.036\,00$, $5\times10^{-7}\%$,
    via the Chern--Simons correction to the QED loop).

  \item \textbf{[Resolved in Paper~IV~\cite{Paper4}]}
    Derive the UV-to-golden-angle running $\Delta_1 = 12.824$ from
    the WZW beta function, identifying the physical scale ratio
    $\mu_{\mathrm{UV}}/\mu_{\mathrm{IR}}$ in condensate units.
    The Compton-wavelength identification and one-loop mass renormalization
    of Paper~IV provide this derivation (Theorems~3.1 and~4.3 therein).

  \item \textbf{[Resolved in Paper~IV~\cite{Paper4}]}
    Determine the thermal correction to the $\alpha$ formula.
    The three- and four-term formulas of Paper~IV are unconditional;
    the four-term formula achieves $5\times10^{-7}\%$. The QED running between $m_{\mathrm{WZW}}$ and $m_e$;
    the thermal correction $(1/\vp)\log(m_e/T_{\mathrm{CMB}})$ is
    subsumed into the spoke formula for $m_e/\Lambda_{\mathrm{cond}}$.
\end{enumerate}

%══════════════════════════════════════════════════════════════════════════════
\section{Summary}
\label{P3:sec:summary}
%══════════════════════════════════════════════════════════════════════════════

We have established a complete analytic route from the icosahedral
Hopfion condensate to two Standard Model constants.

\paragraph{The fine structure constant.}
The cascade $\vp^6/\sin^4(\pi/5) \to 360/\vp^2 \to \aInv$
identifies $\alpha$ as a two-step geometric flow: WZW anti-screening
from the condensate UV scale to the golden-angle attractor, followed by
the one-loop correction $k/(2\pi) = 3/(2\pi)$.  The formula
$\aInv = 360/\vp^2 - k/(2\pi) = 137.030$ agrees with the measured
$137.036$ to $0.004\%$ with no fitted parameters.
The residual $0.006$ is a genuine formula gap, not a lattice artefact
(Section~\ref{P3:sec:lsc_alpha_link}).

\paragraph{The Linking Scale Conjecture.}
Lemma~\ref{P3:lem:bridge} ($\mu^* = \sqrt{k+2}/\vp$, exact) and
Theorem~\ref{P3:thm:lsc_conditional} reduce the LSC to the virial condition
$V^*=\vp$.  The three-fermion bosonization (Section~\ref{P3:sec:bosonization})
identifies the density-feedback model as the $\mathrm{SU}(2)_3$ WZW model
of $k=3$ free Dirac fermions.  The $Q=2$ Hopfion is the $j=0$ singlet;
the Verlinde formula gives $V^* = S_{0,1/2}/S_{0,0} = \vp$ exactly in the
thin-torus.  Theorem~\ref{P3:thm:V_exact_all_R0} (Brouwer + Rellich--Kondrachov)
extends this to all finite $R_0$, closing the LSC completely.

\paragraph{What remains.}
One open item remains from this paper.
The Hessian spectral gap condition (Section~\ref{P3:sec:open}) would
give uniqueness of the self-consistent state; it is not required for any
proved result and does not affect the LSC or $\alpha$ formula.
The genuine $\alpha$ residual --- the formula $360/\vp^2 - k/(2\pi)
= 137.030$ being $0.006$ below the measured $137.036$ --- is resolved
in Paper~IV~\cite{Paper4}: the three-term formula
($\aInv = 137.036\,49$, $0.0004\%$, via the $T$-matrix correction
$T_{1/2}/Q = 3/400$ of the $j=1/2$ spoke carrier)
and the four-term formula ($\aInv = 137.036\,00$, $5\times10^{-7}\%$,
via the Chern--Simons correction to the QED loop).

\begin{center}
\renewcommand{\arraystretch}{1.4}
\begin{tabular}{llcp{3.5cm}}
\toprule
Result & Status & Value & Gap \\
\midrule
$\aInv_{\mathrm{geo}} = \vp^6/\sin^4(\pi/5)$ & Proved & $150.332$ & $9.7\%$ \\
Lemma: $\mu^* = \sqrt{k+2}/\vp$ & \textcolor{darkgreen}{Proved} & exact & --- \\
Lemma: $1+\mu^*\vp = 2\vp$ & \textcolor{darkgreen}{Proved} & exact & --- \\
LSC given $V^*=\vp$ (Thm~\ref{P3:thm:lsc_conditional}) & \textcolor{darkgreen}{Proved} & conditional & --- \\
$V^*=\vp \Leftrightarrow$ LSC $\Leftrightarrow$ $\leff^*=\lam/2^{1/3}$ (Thm~\ref{P3:thm:lsc_equivalence}) & \textcolor{darkgreen}{Proved} & exact & --- \\
Scale invariance: $J_{2a},J_4$ C-indep. & \textcolor{darkgreen}{Proved} & exact & --- \\
EL condition: $d/db[\sqrt{b}I(b)]=0$ & \textcolor{darkgreen}{Proved} & $b^*\!\approx 0.78$ & --- \\
$V(0;R_0)\to\vp$ as $R_0\to\infty$ & \textcolor{darkgreen}{Proved} & $1/R_0^2$ rate & --- \\
$1+1/\vp=\vp$ [key identity] & Exact & $\vp$ & --- \\
Golden angle $= 360/\vp^2$ & Identity & $137.508^\circ$ & --- \\
$\aInv = 360/\vp^2 - k/(2\pi)$ & Result & $137.030$ & $0.004\%$ (improved to $5\times10^{-7}\%$ in Paper~IV~\cite{Paper4}) \\
Exact virial: $V(b)=(15/8)\arctan(\sqrt{8b})/\sqrt{8b}$ & \textcolor{darkgreen}{Proved} & exact & --- \\
$\kappa_{\mathrm{tube}} = g_{\mathrm{WZW}}$ (BS identity) & \textcolor{darkgreen}{Proved} & exact & --- \\
$Q=2$ Hopfion $= j{=}0$ WZW singlet & \textcolor{darkgreen}{Proved} & $h_0{<}h_1$ & --- \\
$V^*=\vp$ via Verlinde (thin-torus) & \textcolor{darkgreen}{Proved} & exact & --- \\
$J_4/J_{2a}=2^{4/3}/\vp^5$ (all $R_0$, Thm~\ref{P3:thm:lsc_1d_complete}) & \textcolor{darkgreen}{Proved} & exact & --- \\
$V^*=\vp$ exactly (all $R_0$), Thm~\ref{P3:thm:V_exact_all_R0} & \textcolor{darkgreen}{Proved} & exact & lattice only \\
Hessian spectral gap (uniqueness) & \textcolor{darkorange}{Open} & --- & --- \\
\bottomrule
\end{tabular}
\end{center}

%══════════════════════════════════════════════════════════════════════════════
\begin{thebibliography}{99}
%══════════════════════════════════════════════════════════════════════════════

\bibitem{Paper1}
F. Manfredi,
\textit{The Density-Feedback Faddeev--Niemi Hopfion:
  Exact Algebraic Predictions for Standard-Model Parameters},
  Zenodo (2026)
\href{https://doi.org/10.5281/zenodo.19342027}{doi:10.5281/zenodo.19342027}

\bibitem{Paper2}
F. Manfredi,
\textit{BPS Structure and Fixed-Point Theorem for the
  Density-Feedback Faddeev--Niemi Hopfion},
\href{https://doi.org/10.5281/zenodo.19392418}{doi:10.5281/zenodo.19392418}

\bibitem{Paper4}
F.~Manfredi,
\textit{The Hopf Spoke, WZW Fermion Mass Renormalization,
  and Three- and Four-Term Formulas for the Fine Structure Constant},
Zenodo (2026).
\href{https://doi.org/10.5281/zenodo.19504729}{doi:10.5281/zenodo.19504729}

\bibitem{PDG2024}
R.~L. Workman et al.\ (Particle Data Group),
\textit{Review of Particle Physics},
Prog.\ Theor.\ Exp.\ Phys.\ \textbf{2022}, 083C01 (2022);
updated online 2024.
\href{https://doi.org/10.1093/ptep/ptac097}{doi:10.1093/ptep/ptac097}

\bibitem{KZ1984}
V.~G. Knizhnik and A.~B. Zamolodchikov,
\textit{Current algebra and Wess--Zumino model in two dimensions},
Nucl.\ Phys.\ B \textbf{247}, 83--103 (1984).
\href{https://doi.org/10.1016/0550-3213(84)90374-2}{doi:10.1016/0550-3213(84)90374-2}

\bibitem{Verlinde1988}
E.~Verlinde,
\textit{Fusion rules and modular transformations in 2D conformal
  field theory},
Nucl.\ Phys.\ B \textbf{300}, 360--376 (1988).
\href{https://doi.org/10.1016/0550-3213(88)90603-7}{doi:10.1016/0550-3213(88)90603-7}

\bibitem{Cardy1986}
J.~L.~Cardy,
\textit{Operator content of two-dimensional conformally invariant theories},
Nucl.\ Phys.\ B \textbf{270}, 186--204 (1986).
\href{https://doi.org/10.1016/0550-3213(86)90552-3}{doi:10.1016/0550-3213(86)90552-3}

\bibitem{Witten1984}
E.~Witten,
\textit{Non-abelian bosonization in two dimensions},
Commun.\ Math.\ Phys.\ \textbf{92}, 455--472 (1984).
\href{https://doi.org/10.1007/BF01215276}{doi:10.1007/BF01215276}

\bibitem{Evans2010}
L.~C.~Evans,
\textit{Partial Differential Equations}, 2nd ed.,
Graduate Studies in Mathematics \textbf{19},
American Mathematical Society (2010).
\href{https://doi.org/10.1090/gsm/019}{doi:10.1090/gsm/019}

\bibitem{Hardy1949}
G.~H.~Hardy,
\textit{Divergent Series},
Oxford University Press, Oxford (1949); repr.\ AMS Chelsea Publishing (1991);
ISBN 978-1-4704-7785-1.


\bibitem{Zumino1984}
B.~Zumino,
\textit{Chiral anomalies and differential geometry},
in: Relativity, Groups and Topology II, Les~Houches (1983),
eds.\ B.S.~DeWitt and R.~Stora, North-Holland, Amsterdam (1984), 1291.

\bibitem{Faddeev1984}
L.~D. Faddeev,
\textit{Operator anomaly for the photon},
Phys.\ Lett.\ B \textbf{145}, 81 (1984).
\href{https://doi.org/10.1016/0370-2693(84)90952-3}{doi:10.1016/0370-2693(84)90952-3}


\bibitem{Paper7}
F.~Manfredi,
\textit{Gravity, Inflation, and Cosmology from the Density-Feedback Hopfion},
Preprint (2026).

\end{thebibliography}


%══════════════════════════════════════════════════════════════════════════════
\appendix
%══════════════════════════════════════════════════════════════════════════════

\section*{Appendix: Pohozaev methods and structural analysis}
\addcontentsline{toc}{section}{Appendix: Pohozaev methods and structural analysis}

\emph{This appendix records the Pohozaev integration-by-parts results and
structural analysis of the 2D proof problem developed during the preparation
of this paper, together with a fourth Pohozaev variation (radial dilation)
added subsequently.  These results establish useful integral identities
that support the bosonization argument in Section~\ref{P3:sec:bosonization},
but are not needed for the main proof chain.
The radial variation (d) connects to the open conjecture
$\vp^9\sqrt{\Ja\Jfour}=Q$ of Paper~IV~\cite{Paper4} and establishes
that conjecture as a bulk--boundary resonance condition.}

%══════════════════════════════════════════════════════════════════════════════
\section{Structural analysis of the 2D proof problem}
\label{P3:sec:2d_structure}

This section provides a complete decomposition of what is needed to
prove any one of the three equivalent conditions of
Theorem~\ref{P3:thm:lsc_equivalence} analytically in the full 2D setting.

\subsection{The key algebraic identity: $\chi = f_{\mathrm{DG}}$}

Define the tube-radial derivative field
\begin{equation}
  \chi(r,z) \;:=\; (r-R_0)\,f_r + z\,f_z \;=\; f_{\mathrm{DG}},
  \label{P3:eq:chi_def}
\end{equation}
which is precisely the quantity $f_{\mathrm{DG}}$ appearing in the
$F_{13}$ component of the solver: $F_{13} = -\sin^2(f)/D_0\cdot\chi$.

\begin{lemma}[BS profile identity]
\label{P3:lem:chi_sinf}
For the Battye--Sutcliffe profile $f = 2\arctan(C/\rho)$,
$\rho = \sqrt{(r-R_0)^2+z^2}$, one has
\begin{equation}
  \chi \;=\; \rho\,f'(\rho) \;=\; -\sin f
  \label{P3:eq:chi_sinf}
\end{equation}
exactly.
\end{lemma}

\begin{proof}
For $f=2\arctan(C/\rho)$: $f'(\rho) = -2C/(\rho^2+C^2)$, so
$\rho f'(\rho) = -2C\rho/(\rho^2+C^2) = -\sin f$.
\end{proof}

This identity is the algebraic key to the 2D proof: it says that the
natural Pohozaev vector $\chi$ coincides (for the BS profile) with
the natural amplitude-scaling direction $-\sin f$.

\subsection{The four Pohozaev variations and what they give}

At the EL saddle of $E_{\mathrm{fb}} = K_{\mathrm{fb}}\cdot J_4$,
the identity $\int \xi\cdot(\delta E_{\mathrm{fb}}/\delta f)\,dV = 0$
holds for any test function $\xi$.  Four natural choices each yield
a distinct integral identity.  Their current status, in light of the
bosonization proof, is recorded here.

\paragraph{Variation (a): Tube-width scaling.}
$\xi_a = \mathrm{d}f_{C}/\mathrm{d}C\big|_{C=C^*}$ for the BS family
$f_C(\rho) = 2\arctan(C/\rho)$.  This gives
$\mathrm{d}E_{\mathrm{fb}}[f_C]/\mathrm{d}C\big|_{C^*} = 0$,
equivalent to $\mathrm{d}/\mathrm{d}b[\sqrt{b}\,I(b)]=0$,
the thin-torus EL condition proved in~\cite{Paper2}.
\emph{Status: proved for thin-torus; does not extend to full 2D
without additional structure.  Superseded by the bosonization proof.}

\paragraph{Variation (b): Tube-expansion scaling.}
$\xi_b = \chi = (r-R_0)f_r + z f_z$.  This gives the Pohozaev
virial identity
\begin{equation}
  J_{2a}\,N_a + \mu^*\,J_{\mathrm{fb}}\,N_{\mathrm{fb}}
  + K_{\mathrm{fb}}\,N_4 = 0,
  \label{P3:eq:pohozaev_b}
\end{equation}
where $N_a,N_4,N_{\mathrm{fb}}$ are the tube-scaling exponents,
and rearranges to
\begin{equation}
  V^* = -\frac{N_a + N_4}{\mu^*\,(N_{\mathrm{fb}} + N_4)}.
  \label{P3:eq:virial_V}
\end{equation}
This identity is exact at the EL saddle (confirmed numerically:
$E_{\mathrm{fb}}[f_\lambda]$ has a local minimum at $\lambda=1$).
It expresses $V^*$ in terms of the scaling exponents, but cannot
determine $V^*=\vp$ without knowing the exponent ratio
$-(N_a+N_4)/(N_{\mathrm{fb}}+N_4) = \sqrt{5}$.
\emph{Status: identity exact and proved; selecting $V^*=\vp$ from it
requires additional input, which the bosonization supplies.}

\paragraph{Variation (c): Amplitude scaling.}
$\xi_c = \sin f$.  By Lemma~\ref{P3:lem:chi_sinf}, this coincides with
$-\chi$ for the BS profile, so variation~(c) is the thin-torus limit
of variation~(b), differing in 2D by $O(\rho/R_0)$.

The amplitude Pohozaev integrals $P_G := \int\!\sin f\cdot(\delta G/\delta f)\,dV$
evaluate exactly by integration by parts (Proposition~\ref{P3:prop:pohozaev_a}),
reducing the identity to a single algebraic condition~\eqref{P3:eq:star_identity}.
This condition is equivalent to $V^*=\vp$ (Proposition~\ref{P3:prop:poh_a_equiv}),
but cannot establish it: the identity holds at the EL saddle for any
value of $V$, not specifically for $\vp$.
\emph{Status: identities proved exactly (Proposition~\ref{P3:prop:pohozaev_a});
cannot select $V^*=\vp$ without the WZW input now supplied by
Section~\ref{P3:sec:bosonization}.}

\paragraph{Variation (d): Radial dilation.}
\label{P3:sec:pohozaev_d}
$\xi_d = r\,\partial_r f$ (the generator of $r\mapsto r(1+\varepsilon)$,
$z$ and $R_0$ fixed).
Unlike variations (a)--(c), this breaks the $(r,z)$ symmetry and
generates a \emph{topological boundary term} $\mathcal{B}(R_0)$ from
the inward shift of the torus core.
The identity $P_{\Kfb}=2\vp P_{\Ja}$ is an algebraic consequence
of $\vp^2=\vp+1$ (proved in Paper~I~\cite{Paper1},
Proposition~2.5).
Substituting the Derrick balance and $P_{\Ja}=-\Ja$ (exact in the
thin-torus: the core shifts to $r=R_0/(1+\varepsilon)$, giving
$J_{2a,\varepsilon}=(2\pi)^2 R_0/(1+\varepsilon)\cdot I_{2a}$)
into the second Pohozaev identity $\Jfour P_{\Kfb}+\Kfb P_{\Jfour}=0$
gives a linear equation for $\Ja$:
\begin{equation}
  \Ja = \frac{\vp^5\,\mathcal{B}(R_0)}{2\,(1+2^{4/3})},
  \label{P3:eq:Ja_from_poh_d}
\end{equation}
where $\mathcal{B}(R_0)=\oint_{\partial B_\varepsilon}F_{13}\cdot(r-R_0)\rho^{-1}dS$
is the Hopf flux of $F_{13}$ through the core tube.
The conjecture $\vp^9\sqrt{\Ja\Jfour}=Q$ holds if and only if
\begin{equation}
  \mathcal{B}(R_0=3) = \frac{2Q\,(1+2^{4/3})}{\vp^{23/2}\cdot 2^{2/3}}
  \approx 0.175.
  \label{P3:eq:B_required}
\end{equation}
The thin-torus formula $4\pi^2 Q_{\mathrm{math}}/R_0\approx26.3$
at $R_0=3$ fails by a factor $\approx150$ because the condensate
at $R_0=3$ is a thick torus ($C^*/R_0\approx0.6$).
\emph{Status: the $P_{\Kfb}=2\vp P_{\Ja}$ identity is proved exactly;
the conjecture reduces to computing $\mathcal{B}(R_0=3)$ from the
full 2D EL profile and showing it equals~\eqref{P3:eq:B_required}.
This single remaining step is equivalent to the original conjecture.
See Open Problem below.}

\begin{remark}[The conjecture as a bulk--boundary resonance condition]
\label{P3:rem:resonance}
The Whitehead/linking-integral structure explains why no local formula
can close the open step above.
For the torus-knot ansatz
$\mathbf{n}=(\sin f\cos Q_{\mathrm{math}}\phi,\sin f\sin Q_{\mathrm{math}}\phi,\cos f)$
the Hopf curvature is $F=Q_{\mathrm{math}}\sin f\,df\wedge d\phi$
(no $dr\wedge dz$ component), the connection
$A=Q_{\mathrm{math}}(1-\cos f)\,d\phi$ satisfies $dA=F$, and
\begin{equation}
  A\wedge F = Q_{\mathrm{math}}^2(1-\cos f)\sin f\,(d\phi\wedge df\wedge d\phi) = 0,
\end{equation}
so the Whitehead integral vanishes identically in the bulk.
The Hopf charge $Q_{\mathrm{math}}$ lives entirely in the seam/transition
function $\lambda_{NS}=Q_{\mathrm{math}}\phi$ on $S^3$:
$Q_{\mathrm{math}}=(1/2\pi)\int_{\mathrm{seam}}(A_N-A_S)=Q_{\mathrm{math}}$
for \emph{any} continuous torus-knot field, independent of $f(r,z)$.

The conjecture
$\vp^9\sqrt{\Ja\Jfour}=Q_{\mathrm{paper}}=10$
therefore equates a continuous bulk functional (left side, depends on
$f(r,z)$ everywhere) to a discrete group-theoretic integer $\mathrm{ord}(q_{2I})$
(right side, independent of $f(r,z)$).
No Pohozaev identity, BPS bound, Whitehead integral, or seam computation
can establish this: all such tools either give identities holding for
every smooth profile, or give the topological integer
$Q_{\mathrm{math}}=2$ which constrains only the seam, not the bulk shape.
The conjecture is a \emph{resonance condition}: the EL solution at the
specific condensate parameters $\lam=\vp^6$, $\ms=3-\vp$, $R_0=3$
is tuned so that a bulk integral exactly equals a group-theoretic integer.
Only the EL equation itself can enforce this.
\end{remark}

Before the bosonization proof (Section~\ref{P3:sec:bosonization}), the
proof of $V^*=\vp$ was decomposed into three sub-problems.  Their
current status:

\begin{description}
\item[Sub-problem A (Amplitude Pohozaev identity) --- partial.]
  The amplitude Pohozaev identity $\int \sin f \cdot (\delta E_{\mathrm{fb}}/\delta f)\,dV = 0$
  at the EL saddle is exact (proved in Section~\ref{P3:sec:pohozaev_a}).
  However, it gives only that \emph{some} $V$ satisfies the identity
  at the saddle; it cannot distinguish $V=\vp$ from other values without
  additional WZW input.  Sub-problem~A established useful integral identities
  (Proposition~\ref{P3:prop:pohozaev_a}) but could not close the proof alone.

\item[Sub-problem B (Tube-scaling minimum) --- open, not needed.]
  Prove analytically that the EL saddle is a local minimum of
  $E_{\mathrm{fb}}[f_\lambda]$ under tube scaling at $\lambda=1$.
  Numerically confirmed ($E_{\mathrm{fb}}$ increases for $\lambda\neq 1$;
  see Section~\ref{P3:sec:saddle_num}).  This is the Hessian spectral gap
  condition (Open Problem~1 of Section~\ref{P3:sec:open}), and is not
  required for the main results.

\item[Sub-problem C (WZW selection of $\mu^*$) --- \textbf{closed}.]
  $\mu^* = \sqrt{k+2}/\vp$ is established by Lemma~\ref{P3:lem:bridge}
  (algebraic) and given its physical origin by the three-fermion
  bosonization (Section~\ref{P3:sec:bosonization}): $\mu^*$ is the
  Wess--Zumino coupling constant of the $k=3$ free-fermion system,
  and the specific value $\sqrt{k+2}/\vp$ is the unique WZ coupling
  that makes the model equivalent to $\mathrm{SU}(2)_3$ WZW.
\end{description}

The closure route that succeeded was not sub-problem~A (Pohozaev) but
sub-problem~C (WZW identification), combined with the Verlinde formula
to give $V^* = S_{0,1/2}/S_{0,0} = \vp$ directly.

\subsection{Numerical confirmation of the saddle structure}
\label{P3:sec:saddle_num}

From the converged $R_0=3$ profile:

\begin{center}
\renewcommand{\arraystretch}{1.3}
\begin{tabular}{lcc}
\toprule
Quantity & Value & WZW prediction \\
\midrule
$V^* = \Jfb/\Ja$ & $1.61824$ & $\vp = 1.61803$ \\
$J_4/J_{2a}$ & $0.22717$ & $2^{4/3}/\vp^5 = 0.22721$ \\
$\leff^* = K_{\mathrm{fb}}/J_4$ & $14.2465$ & $\vp^6/2^{1/3} = 14.2424$ \\
$E_{\mathrm{fb}}(f_\lambda)$ at $\lambda=0.98$ & $1.131\times E_0$ & --- \\
$E_{\mathrm{fb}}(f_\lambda)$ at $\lambda=1.02$ & $1.099\times E_0$ & --- \\
\bottomrule
\end{tabular}
\end{center}

The last two rows confirm Sub-problem~B numerically: $E_{\mathrm{fb}}$
increases in both scaling directions, confirming the EL is a
tube-scaling local minimum.  The gaps in $V^*$, $J_4/J_{2a}$, and $\leff^*$
are $0.013$--$0.029\%$, consistent with lattice artefacts at $h=0.12$.
They do not appear in the continuum: $V^*=\vp$ is proved exactly
by Theorem~\ref{P3:thm:V_exact_all_R0}.


%══════════════════════════════════════════════════════════════════════════════

%══════════════════════════════════════════════════════════════════════════════
\section{The amplitude Pohozaev identity}
\label{P3:sec:pohozaev_a}
%══════════════════════════════════════════════════════════════════════════════

\subsection{Exact evaluation of the amplitude Pohozaev integrals}

\begin{proposition}[Amplitude Pohozaev integrals]
\label{P3:prop:pohozaev_a}
Let $f$ be any smooth $Q=2$ profile.  Define the cos-weighted functionals
\begin{equation}
  J_{2a}^{cf} := \int \cos f\cdot\sin^4\!f\cdot\kk\,dV,\quad
  J_4^{cf}    := \int \cos f\cdot (F_{13}^2+F_{12}^2+F_{23}^2)\,dV,\quad
  J_{\mathrm{fb}}^{cf2} := \int \cos f\cdot \kk\cdot f_\beta^2\,dV,
\end{equation}
where $f_\beta = 1/(1+\beta\kk)$.  Then
\begin{align}
  \int \sin f\cdot\frac{\delta J_{2a}}{\delta f}\,dV
    &= 6\,J_{2a}^{cf},      \label{P3:eq:poh_J2a}\\
  \int \sin f\cdot\frac{\delta J_{\mathrm{fb}}}{\delta f}\,dV
    &= 2\,J_{\mathrm{fb}}^{cf2}, \label{P3:eq:poh_Jfb}\\
  \int \sin f\cdot\frac{\delta J_4}{\delta f}\,dV
    &= 6\,J_4^{cf}.          \label{P3:eq:poh_J4}
\end{align}
\end{proposition}

\begin{proof}
For each functional, we split $\delta G/\delta f$ into a local part
(depending on $f$ but not on $\nabla f$) and a divergence part (from
$\nabla f$), then apply the divergence theorem.

\textbf{Proof of~\eqref{P3:eq:poh_J2a}.}
The local part of $\delta J_{2a}/\delta f$ is
$4\sin^3\!f\cos f\cdot\kk + 2\sin^5\!f\cos f\cdot A$,
giving $\int\!\sin f\cdot\text{local} = 4J_{2a}^{cf} + 2\int\!\sin^6\!f\cos f\cdot A\,dV$.
The flux is $2\sin^4\!f\,\nabla f$; by IBP,
$\int\!\sin f\cdot(-\mathrm{div}(2\sin^4\!f\,\nabla f))\,dV
= 2\int\!\cos f\,\sin^4\!f\,|\nabla f|^2\,dV$.
Summing and using $\kk = |\nabla f|^2 + \sin^2\!f\cdot A$:
$6J_{2a}^{cf}$.

\textbf{Proof of~\eqref{P3:eq:poh_Jfb}.}
Local part: $2\sin f\cos f\cdot A\cdot f_\beta^2$.
Flux: $2f_\beta^2\,\nabla f$.  IBP gives $2\int\!\cos f\,f_\beta^2|\nabla f|^2\,dV$.
Sum: $2\int\!\cos f\,f_\beta^2\kk\,dV = 2J_{\mathrm{fb}}^{cf2}$.

\textbf{Proof of~\eqref{P3:eq:poh_J4}.}
Using $F_{13}=-\sin^2\!f/D_0\cdot\chi$ and $F_{12}=\sin^2\!f\,f_r/r$,
$F_{23}=\sin^2\!f\,f_z/r$: the local part gives $4J_4^{cf}$;
IBP of the gradient fluxes gives $2J_4^{cf}$.  Total: $6J_4^{cf}$.
\end{proof}

\begin{remark}
The factor 6 in~\eqref{P3:eq:poh_J2a} and~\eqref{P3:eq:poh_J4} vs.\ factor 2
in~\eqref{P3:eq:poh_Jfb} reflects the different feedback structure:
$J_{2a}$ and $J_4$ have $\sin^4\!f$ or $\sin^2\!f$ prefactors (both
giving IBP contributions of $+2$), while $J_{\mathrm{fb}}$ has the
nonlinear feedback weight $f_\beta^2$ which does not contribute an
additional factor.
\end{remark}

\subsection{The amplitude Pohozaev identity at the EL saddle}

At the EL saddle of $E_{\mathrm{fb}} = \Kfb\cdot\Jfour$,
the Pohozaev identity with $\xi = \sin f$ gives
\begin{equation}
  \Jfour\cdot P_{\Kfb} + \Kfb\cdot P_{\Jfour} = 0
  \label{P3:eq:pohozaev_efb}
\end{equation}
where $P_G := \int\!\sin f\cdot(\delta G/\delta f)\,dV$.
Substituting Proposition~\ref{P3:prop:pohozaev_a}:
\begin{equation}
  \Jfour\cdot(6\,J_{2a}^{cf} + 2\mu^*\,J_{\mathrm{fb}}^{cf2})
  + \Kfb\cdot 6\,J_4^{cf} = 0.
  \label{P3:eq:poh_A_full}
\end{equation}
Setting $a := J_{2a}^{cf}/\Ja$, $b := J_4^{cf}/\Jfour$,
$c := J_{\mathrm{fb}}^{cf2}/\Jfb$ and $V := \Jfb/\Ja$,
equation~\eqref{P3:eq:poh_A_full} yields:
\begin{equation}
  V^* \;=\; \frac{-3(a+b)}{\mu^*\,(c+3b)}.
  \label{P3:eq:V_from_poh_A}
\end{equation}

\subsection{Reduction to a single algebraic identity}

\begin{proposition}[Main identity]
\label{P3:prop:star_identity}
$V^* = \vp$ if and only if the following identity holds at the EL saddle:
\begin{equation}
  \sqrt{k+2}\,\frac{J_{\mathrm{fb}}^{cf2}}{\Jfb}
  \;+\; 3\!\left(\frac{J_{2a}^{cf}}{\Ja} + 2\vp\,\frac{J_4^{cf}}{\Jfour}\right)
  \;=\; 0.
  \label{P3:eq:star_identity}
\end{equation}
With $k=3$: $\sqrt{k+2} = \sqrt{5} = \mu^*\vp$ (Lemma~\ref{P3:lem:bridge}).
\end{proposition}

\begin{proof}
From~\eqref{P3:eq:V_from_poh_A}, $V^* = \vp$ iff
$\mu^*\vp(c+3b) = -3(a+b)$, i.e.\
$\sqrt{5}(c+3b)+3(a+b)=0$, i.e.\
$\sqrt{5}\,c+3(a+(1+\sqrt{5})b)=0$, i.e.\
$\sqrt{5}\,c+3(a+2\vp b)=0$.
Multiplying through by $(J_{\mathrm{fb}}^{cf2}/J_{\mathrm{fb}}^{cf2})^{-1}$ etc.\ gives~\eqref{P3:eq:star_identity}.
\end{proof}

\begin{remark}[Physical interpretation]
Identity~\eqref{P3:eq:star_identity} relates the \emph{cos-weighted} averages of
the three curvature functionals.  The cos-weight $\cos f$ is the angular factor
in the spin-connection of the Hopfion field $n: \mathbb{R}^3\to S^2$, and the
identity is a constraint on how the profile distributes curvature between
tube, winding, and feedback channels at the WZW-selected $\mu^*$.
\end{remark}

\subsection{Numerical verification}

From the converged $R_0=3$, $\bst=0.452$ profile:
\begin{center}
\renewcommand{\arraystretch}{1.3}
\begin{tabular}{lcc}
\toprule
Quantity & Value & Formula result \\
\midrule
$a = J_{2a}^{cf}/J_{2a}$ & $0.06483$ & --- \\
$b = J_4^{cf}/J_4$ & $-0.07942$ & --- \\
$c = J_{\mathrm{fb}}^{cf2}/J_{\mathrm{fb}}$ & $0.24785$ & --- \\
LHS of~\eqref{P3:eq:star_identity}: $\sqrt{5}\,c + 3(a+2\vp b)$ & $-0.0223$ & $0$ \\
$V$ from~\eqref{P3:eq:V_from_poh_A} & $3.301$ & $\vp = 1.618$ \\
\bottomrule
\end{tabular}
\end{center}
The LHS of~\eqref{P3:eq:star_identity} evaluates to $-0.022$ (not $0$),
a $4\%$ gap relative to $\sqrt{5}|c|$.  The Pohozaev-derived
$V = 3.301$ differs substantially from $\vp = 1.618$.

These results confirm what Remark~\ref{P3:rem:five_forms} states: the
amplitude Pohozaev identity holds for \emph{any} $V$ at the corresponding
EL saddle, not specifically for $V = \vp$.  The identity alone cannot
close the proof; the value $\vp$ is selected by the WZW structure
(Section~\ref{P3:sec:bosonization}), not by the Pohozaev route.
The data above characterise the profile's cos-weighted curvature ratios
$(a, b, c)$, which are useful secondary diagnostics of profile convergence.

\begin{proposition}[Equivalence of Pohozaev-A form]
\label{P3:prop:poh_a_equiv}
Identity~\eqref{P3:eq:star_identity} is equivalent to $V^* = \vp$ at the EL saddle:
\begin{equation}
  \sqrt{k+2}\,\frac{J_{\mathrm{fb}}^{cf2}}{\Jfb} + 3\!\left(\frac{J_{2a}^{cf}}{\Ja}
  + 2\vp\,\frac{J_4^{cf}}{\Jfour}\right) = 0
  \;\;\Longleftrightarrow\;\; V^* = \vp.
\end{equation}
\end{proposition}

\begin{proof}
The amplitude Pohozaev identity at the EL saddle is
$3(a+b) + \mu^* V(3b+c) = 0$ (from~\eqref{P3:eq:poh_A_full}).
Substituting $V = \vp$ and $\mu^*\vp = \sqrt{5} = \sqrt{k+2}$ (Lemma~\ref{P3:lem:bridge})
gives~\eqref{P3:eq:star_identity} directly.
Conversely, given~\eqref{P3:eq:star_identity}, the same identity gives $V=\vp$.
\end{proof}

\begin{remark}[Equivalent forms, now all proved]
\label{P3:rem:five_forms}
The following conditions are all equivalent at the EL saddle
(by Theorem~\ref{P3:thm:lsc_equivalence} and Proposition~\ref{P3:prop:poh_a_equiv}),
and all proved via the bosonization (Section~\ref{P3:sec:bosonization}):
\begin{enumerate}[noitemsep,label=(\roman*)]
  \item $V^* = \vp$ \hfill[virial form --- proved]
  \item $J_4/J_{2a} = 2^{4/3}/\vp^5$ \hfill[LSC form --- proved]
  \item $\leff^* = \vp^6/2^{1/3}$ \hfill[scale form --- proved]
  \item $\sqrt{k+2}\,c + 3(a+2\vp b) = 0$ \hfill[Pohozaev-A form --- proved]
\end{enumerate}
The numerical values confirm (i)--(iii) to within lattice artefacts of
$0.013$--$0.029\%$.  Form~(iv) is the concrete integral identity of
this appendix; it holds when $V^*=\vp$, which is proved in
Theorem~\ref{P3:thm:V_exact_all_R0}.

The Pohozaev identity $3(a+b)+\mu^* V(3b+c)=0$ holds for any value
of $V$ at the EL saddle; it characterises the saddle structure but
cannot by itself select $V=\vp$.  The selection is made by the WZW
quantum dimension $\dim_q(\tfrac{1}{2};k{=}3) = \vp$ via bosonization.
\end{remark}


%══════════════════════════════════════════════════════════════════════════════

\end{document}